\documentclass[fs13,seriesSau]{korfarm-base}
\usepackage{korfarm-saussure}

% ============================================================
% passage-note.tex (v11)
% 지문 박스 + 첫 페이지 우측 메모란 (동적 높이)
% 정책: PAGE_BREAK_POLICY.md §3, §4 참조
%
% 변경 (v11):
%  · 지문 박스 폭 확대: enlarge right by=-3.7cm  (메모 3.4cm + gap 3mm)
%  · 메모란 동적 높이: frame.north east ~ frame.south east 사이 자동 채움
%  · 메모 줄선: clip 영역으로 박스 내부만 채움 (박스 높이 모르고도 작동)
%  · 문단 들여쓰기 활성화: tcolorbox 내부 \parindent=1em, \noindent 제거
%
% 사용:
%   \kfPassageNote{제목}{지문 본문}      → 메모란 포함
%   \kfPassageNoteNT{지문 본문}          → 메모란 포함, 제목 없음
%   \kfPassageFull{제목}{본문}           → 메모란 없음 (활동 내 보기)
% ============================================================

\makeatletter

% 메모란 규격 (cm 고정)
\newlength{\kf@memoW}
\setlength{\kf@memoW}{3.4cm}
\newlength{\kf@memoGap}
\setlength{\kf@memoGap}{3mm}

% --- 메인: 제목 있는 버전 ---
\newcommand{\kf@passagenote@with}[2]{%
  \par\ifkfPassageNewPage\ifdim\pagetotal>0.4\textheight\clearpage\fi\fi\smallskip\needspace{6\baselineskip}%
  \begin{tcolorbox}[
    enhanced, breakable,
    width=\dimexpr\linewidth-4cm\relax,
    colback=kfPassageBg, colframe=kfPrimary!70,
    boxrule=0.6pt, arc=3pt,
    left=\kfBoxPad, right=\kfBoxPad, top=\dimexpr\kfBoxPad+8pt\relax, bottom=\kfBoxPad,
    title={\small\bfseries\color{kfPrimary} #1},
    coltitle=kfPrimary,
    colbacktitle=kfPrimaryLight!60,
    attach boxed title to top left={yshift=-2.4mm, xshift=8pt},
    boxed title style={colframe=kfPrimary!50, boxrule=0.5pt, arc=2pt},
    parbox=false,
    overlay unbroken and first={%
      \begin{scope}
        \draw[kfRule, line width=0.4pt, fill=kfNoteBg, rounded corners=2pt]
          ([xshift=\kf@memoGap]frame.north east) rectangle
          ([xshift=\kf@memoGap+\kf@memoW]frame.south east);
        \node[anchor=north east, font=\footnotesize\bfseries\color{kfMute}, inner sep=3pt]
          at ([xshift=\kf@memoGap+\kf@memoW, yshift=-2pt]frame.north east) {메모};
        \node[anchor=north east, fill=kfPrimary!15, text=kfPrimary,
              font=\scriptsize\bfseries, rounded corners=2pt,
              inner xsep=4pt, inner ysep=2pt]
          at ([xshift=\kf@memoGap+\kf@memoW-3pt, yshift=-19pt]frame.north east) {\kf@memocat};
        \begin{scope}
          \path[clip]
            ([xshift=\kf@memoGap]frame.north east) rectangle
            ([xshift=\kf@memoGap+\kf@memoW]frame.south east);
          \foreach \i in {1,...,40}{%
            \draw[kfRule, line width=0.25pt]
              ([xshift=\kf@memoGap+2mm, yshift=-7mm - \i*5mm]frame.north east) --
              ([xshift=\kf@memoGap+\kf@memoW-2mm, yshift=-7mm - \i*5mm]frame.north east);%
          }
        \end{scope}
      \end{scope}
    }
  ]
    \setlength{\parindent}{1em}%
    \hspace*{1em}\ignorespaces#2%
  \end{tcolorbox}\par\smallskip
}

% --- 제목 없는 버전 ---
\newcommand{\kf@passagenote@notitle}[1]{%
  \par\ifkfPassageNewPage\ifdim\pagetotal>0.4\textheight\clearpage\fi\fi\smallskip\needspace{6\baselineskip}%
  \begin{tcolorbox}[
    enhanced, breakable,
    width=\dimexpr\linewidth-4cm\relax,
    colback=kfPassageBg, colframe=kfPrimary!70,
    boxrule=0.6pt, arc=3pt,
    left=\kfBoxPad, right=\kfBoxPad, top=\dimexpr\kfBoxPad+8pt\relax, bottom=\kfBoxPad,
    parbox=false,
    overlay unbroken and first={%
      \begin{scope}
        \draw[kfRule, line width=0.4pt, fill=kfNoteBg, rounded corners=2pt]
          ([xshift=\kf@memoGap]frame.north east) rectangle
          ([xshift=\kf@memoGap+\kf@memoW]frame.south east);
        \node[anchor=north east, font=\footnotesize\bfseries\color{kfMute}, inner sep=3pt]
          at ([xshift=\kf@memoGap+\kf@memoW, yshift=-2pt]frame.north east) {메모};
        \node[anchor=north east, fill=kfPrimary!15, text=kfPrimary,
              font=\scriptsize\bfseries, rounded corners=2pt,
              inner xsep=4pt, inner ysep=2pt]
          at ([xshift=\kf@memoGap+\kf@memoW-3pt, yshift=-19pt]frame.north east) {\kf@memocat};
        \begin{scope}
          \path[clip]
            ([xshift=\kf@memoGap]frame.north east) rectangle
            ([xshift=\kf@memoGap+\kf@memoW]frame.south east);
          \foreach \i in {1,...,40}{%
            \draw[kfRule, line width=0.25pt]
              ([xshift=\kf@memoGap+2mm, yshift=-7mm - \i*5mm]frame.north east) --
              ([xshift=\kf@memoGap+\kf@memoW-2mm, yshift=-7mm - \i*5mm]frame.north east);%
          }
        \end{scope}
      \end{scope}
    }
  ]
    \setlength{\parindent}{1em}%
    \hspace*{1em}\ignorespaces#1%
  \end{tcolorbox}\par\smallskip
}

\newcommand{\kfPassageNote}[2]{%
  \def\kf@tmptitle{#1}%
  \ifx\kf@tmptitle\@empty
    \kf@passagenote@notitle{#2}%
  \else
    \kf@passagenote@with{#1}{#2}%
  \fi
}
\newcommand{\kfPassageNoteNT}[1]{\kf@passagenote@notitle{#1}}
\makeatother

% ============================================================
% 풀폭 지문 박스 (메모란 없음) — 활동 내 보기 박스
% PAGE_BREAK_POLICY.md §3
% ============================================================
\makeatletter
\newcommand{\kf@passagefull@with}[2]{%
  \par\ifkfPassageNewPage\ifdim\pagetotal>0.4\textheight\clearpage\fi\fi\smallskip\needspace{6\baselineskip}%
  \begin{tcolorbox}[
    enhanced, breakable,
    colback=kfPassageBg, colframe=kfPrimary!70,
    boxrule=0.6pt, arc=3pt,
    left=\kfBoxPad, right=\kfBoxPad, top=\dimexpr\kfBoxPad+8pt\relax, bottom=\kfBoxPad,
    title={\small\bfseries\color{kfPrimary} #1},
    coltitle=kfPrimary,
    colbacktitle=kfPrimaryLight!60,
    attach boxed title to top left={yshift=-2.4mm, xshift=8pt},
    boxed title style={colframe=kfPrimary!50, boxrule=0.5pt, arc=2pt},
    parbox=false
  ]
    \setlength{\parindent}{1em}%
    \hspace*{1em}\ignorespaces#2%
  \end{tcolorbox}\par\smallskip
}
\newcommand{\kf@passagefull@notitle}[1]{%
  \par\ifkfPassageNewPage\ifdim\pagetotal>0.4\textheight\clearpage\fi\fi\smallskip\needspace{6\baselineskip}%
  \begin{tcolorbox}[
    enhanced, breakable,
    colback=kfPassageBg, colframe=kfPrimary!70,
    boxrule=0.6pt, arc=3pt,
    left=\kfBoxPad, right=\kfBoxPad, top=\dimexpr\kfBoxPad+8pt\relax, bottom=\kfBoxPad,
    parbox=false
  ]
    \setlength{\parindent}{1em}%
    \hspace*{1em}\ignorespaces#1%
  \end{tcolorbox}\par\smallskip
}
\newcommand{\kfPassageFull}[2]{%
  \def\kf@tmptitle{#1}%
  \ifx\kf@tmptitle\@empty
    \kf@passagefull@notitle{#2}%
  \else
    \kf@passagefull@with{#1}{#2}%
  \fi
}
\newcommand{\kfPassageFullNT}[1]{\kf@passagefull@notitle{#1}}
\makeatother

% ============================================================
% 작은 문장 박스 (문장 독해용) — PAGE_BREAK_POLICY.md §2.5
% ============================================================
\newcommand{\kfSentenceBox}[2]{%
  \par\smallskip\needspace{3\baselineskip}%
  \noindent\begin{tcolorbox}[
    enhanced, breakable=false,
    colback=kfPassageBg, colframe=kfMute,
    boxrule=0.4pt, arc=2pt,
    left=8pt, right=6pt, top=4pt, bottom=4pt,
    title={\footnotesize\bfseries\color{kfPrimary} #1},
    coltitle=kfPrimary, colbacktitle=kfPaper,
    attach boxed title to top left={yshift=-1.6mm, xshift=4pt},
    boxed title style={colframe=kfMute, boxrule=0.3pt, arc=1pt}
  ]%
    \setstretch{1.3}#2%
  \end{tcolorbox}\par\smallskip%
}

% ============================================================
% 지문 본문 아래 안내/정리 박스 (v16 신규)
% 사용: \kfPassageHint{한 줄 정리}{본문 안내 텍스트}
% ============================================================
\newcommand{\kfPassageHint}[2]{%
  \par\smallskip\needspace{4\baselineskip}%
  \begin{tcolorbox}[
    enhanced, breakable=false,
    colback=kfPrimary!5, colframe=kfPrimary!40,
    boxrule=0.4pt, arc=2pt,
    left=8pt, right=6pt, top=4pt, bottom=4pt,
    title={\footnotesize\bfseries\color{kfPrimary} #1},
    coltitle=kfPrimary, colbacktitle=kfPaper,
    attach boxed title to top left={yshift=-1.5mm, xshift=4pt},
    boxed title style={colframe=kfPrimary!40, boxrule=0.3pt, arc=1pt}
  ]%
    \small\setstretch{1.3}#2%
  \end{tcolorbox}\par\smallskip%
}

% ============================================================
% 환경 버전 (호환성) — 메모란 포함
% ============================================================
\makeatletter
\newenvironment{kfPassageNoteEnv}[1]%
  {\par\needspace{6\baselineskip}%
   \def\kf@psrc{#1}%
   \begin{tcolorbox}[
     enhanced, breakable,
     width=\dimexpr\linewidth-4cm\relax,
     colback=kfPassageBg, colframe=kfPrimary!70,
     boxrule=0.6pt, arc=3pt,
     left=\kfBoxPad, right=\kfBoxPad, top=\dimexpr\kfBoxPad+8pt\relax, bottom=\kfBoxPad,
     title={\small\bfseries\color{kfPrimary} \kf@psrc},
     coltitle=kfPrimary, colbacktitle=kfPrimaryLight!60,
     attach boxed title to top left={yshift=-2.4mm, xshift=8pt},
     boxed title style={colframe=kfPrimary!50, boxrule=0.5pt, arc=2pt},
     parbox=false,
     overlay unbroken and first={%
       \begin{scope}
         \draw[kfRule, line width=0.4pt, fill=kfNoteBg, rounded corners=2pt]
           ([xshift=\kf@memoGap]frame.north east) rectangle
           ([xshift=\kf@memoGap+\kf@memoW]frame.south east);
         \node[anchor=north east, font=\footnotesize\bfseries\color{kfMute}, inner sep=3pt]
           at ([xshift=\kf@memoGap+\kf@memoW, yshift=-2pt]frame.north east) {메모};
         \begin{scope}
           \path[clip]
             ([xshift=\kf@memoGap]frame.north east) rectangle
             ([xshift=\kf@memoGap+\kf@memoW]frame.south east);
           \foreach \i in {1,...,40}{%
             \draw[kfRule, line width=0.25pt]
               ([xshift=\kf@memoGap+2mm, yshift=-7mm - \i*5mm]frame.north east) --
               ([xshift=\kf@memoGap+\kf@memoW-2mm, yshift=-7mm - \i*5mm]frame.north east);%
           }
         \end{scope}
       \end{scope}
     }
   ]%
   \setlength{\parindent}{1em}%
  }%
  {\end{tcolorbox}\par\smallskip}
\makeatother

% ============================================================
% condition-box.tex
% <조건> 박스 컴포넌트 (tcolorbox)
% 사용:
%   \begin{kfCondition}
%     \item 첫째 조건
%     \item 둘째 조건
%   \end{kfCondition}
% ============================================================

\newenvironment{kfCondition}%
  {\par\smallskip\needspace{4\baselineskip}%
   \begin{tcolorbox}[
     enhanced, breakable=false,
     colback=kfPrimaryLight!40, colframe=kfPrimary,
     boxrule=0.5pt, arc=2pt,
     left=\kfBoxPad, right=\kfBoxPad, top=\kfBoxPad, bottom=\kfBoxPad,
     title={\small\bfseries\color{kfPrimary} \,조건\,},
     coltitle=kfPrimary, colbacktitle=kfPaper,
     attach boxed title to top left={yshift=-2mm, xshift=6pt},
     boxed title style={colframe=kfPrimary, boxrule=0.4pt, arc=1pt}
   ]%
   \begin{enumerate}[leftmargin=16pt, itemsep=1pt, topsep=1pt,
                    label=\textbf{\arabic*.}]}%
  {\end{enumerate}\end{tcolorbox}\par\smallskip}

% ============================================================
% evidence-box.tex
% <보기> 박스 컴포넌트
% 사용:
%   \begin{kfEvidence}
%     보기 본문 ...
%   \end{kfEvidence}
% ============================================================

\newenvironment{kfEvidence}%
  {\par\smallskip\needspace{4\baselineskip}%
   \begin{tcolorbox}[
     enhanced, breakable=false,
     colback=kfPaper, colframe=kfMute,
     boxrule=0.5pt, arc=2pt,
     left=\kfBoxPad, right=\kfBoxPad, top=\kfBoxPad, bottom=\kfBoxPad,
     title={\small\bfseries\color{kfInk} \,보기\,},
     coltitle=kfInk, colbacktitle=kfPrimaryLight!50,
     attach boxed title to top left={yshift=-2mm, xshift=6pt},
     boxed title style={colframe=kfMute, boxrule=0.4pt, arc=1pt}
   ]%
   \leavevmode\mbox{}\ignorespaces%
  }%
  {\end{tcolorbox}\par\smallskip}

% ============================================================
% choice-list.tex
% 5선지 (① ② ③ ④ ⑤) 자동 들여쓰기 컴포넌트
% 사용:
%   \begin{kfChoices}
%     \kfChoice 첫째 선택지
%     \kfChoice 둘째 선택지
%     \kfChoice 셋째 선택지
%     \kfChoice 넷째 선택지
%     \kfChoice 다섯째 선택지
%   \end{kfChoices}
% ============================================================

\newcounter{kfChoiceCnt}
\newcommand{\kfChoiceMark}[1]{%
  \ifcase#1\relax%
  \or ①\or ②\or ③\or ④\or ⑤\or ⑥\or ⑦\or ⑧\or ⑨%
  \fi
}

% 환경: 자동 번호
\newenvironment{kfChoices}%
  {\setcounter{kfChoiceCnt}{0}%
   \par\smallskip%
   \begin{list}{}{%
     \setlength{\leftmargin}{20pt}%
     \setlength{\itemindent}{0pt}%
     \setlength{\labelwidth}{18pt}%
     \setlength{\labelsep}{6pt}%
     \setlength{\itemsep}{2pt}%
     \setlength{\topsep}{1pt}%
     \setlength{\parsep}{0pt}%
     \renewcommand{\makelabel}[1]{##1\hss}%
   }}%
  {\end{list}\par\smallskip}

\newcommand{\kfChoice}{%
  \stepcounter{kfChoiceCnt}%
  \item[\textbf{\kfChoiceMark{\value{kfChoiceCnt}}}]\ignorespaces%
}

% --- 한 줄 5선지 (짧은 선택지용) ---
\newcommand{\kfChoicesInline}[5]{%
  \par\smallskip%
  \noindent\textbf{①}~#1\quad\textbf{②}~#2\quad\textbf{③}~#3\quad\textbf{④}~#4\quad\textbf{⑤}~#5%
  \par\smallskip%
}

% ============================================================
% answer-note.tex
% 답안 작성란 (노트 스타일, 줄친 영역)
% 사용:
%   \kfAnswerLines{5}        % 5줄 답안란
%   \kfAnswerNote{서술형 답안란}{8} % 라벨+8줄
% ============================================================

% 단일 줄 (필요 시 인라인)
\newcommand{\kfAnswerLine}{%
  \par\noindent\textcolor{kfRule}{\rule{\linewidth}{0.4pt}}\par
}

% N줄 답안란 (간단)
\newcommand{\kfAnswerLines}[1]{%
  \par\smallskip%
  \begin{minipage}{\linewidth}%
  \foreach \i in {1,...,#1}{%
    \textcolor{kfRule}{\rule{\linewidth}{0.4pt}}\\[6pt]%
  }%
  \end{minipage}%
  \par\smallskip%
}

% 라벨이 있는 노트형 답안란
\newcommand{\kfAnswerNote}[2]{%
  \par\smallskip\needspace{#2\baselineskip}%
  \begin{tcolorbox}[
    enhanced, breakable=false,
    colback=kfPaper, colframe=kfRule,
    boxrule=0.4pt, arc=2pt,
    left=\kfBoxPad, right=\kfBoxPad, top=\kfBoxPad, bottom=\kfBoxPad,
    title={\footnotesize\bfseries\color{kfMute} #1},
    coltitle=kfMute, colbacktitle=kfPaper,
    attach boxed title to top left={yshift=-1.5mm, xshift=6pt},
    boxed title style={colframe=kfRule, boxrule=0.3pt, arc=1pt}
  ]%
  \foreach \i in {1,...,#2}{%
    \textcolor{kfRule}{\rule{\linewidth}{0.3pt}}\\[6pt]%
  }%
  \end{tcolorbox}\par\smallskip%
}

% 짧은 빈칸 (단답형)
\newcommand{\kfBlank}[1][3cm]{\underline{\hspace{#1}}}
% 넓은 빈칸 (서술 한 줄)
\newcommand{\kfBlankWide}[1][6cm]{\underline{\hspace{#1}}}
% 괄호 빈칸
\newcommand{\kfBlankParen}{(\,\hspace{1cm}\,)}
% O/X 빈칸
\newcommand{\kfBlankOX}{(\,\hspace{1.5cm}\,)}

% ============================================================
% question-block.tex
% 문제 블록 (번호 배지 + 발문 + 보기/조건 + 선택지 + 답안란)
% 모든 구성을 하나의 minipage로 묶어 단/페이지 단절 금지
%
% 사용:
%   \begin{kfQBlock}{1}{객관식}
%     \kfQStem{다음 글에 대한 설명으로 가장 적절한 것은?}
%     \begin{kfEvidence} ... \end{kfEvidence}    % 선택
%     \begin{kfCondition} ... \end{kfCondition}  % 선택
%     \begin{kfChoices}
%       \kfChoice ...
%     \end{kfChoices}
%     \kfAnswerLines{2}                            % 선택
%   \end{kfQBlock}
% ============================================================

% 무단절 minipage로 묶고, 발문/보기/조건/선택지/답안란을 자유롭게 배치
% 사용자 피드백 #4: [객관식]/[서술형] 등 텍스트 라벨 제거 — 번호 배지만 표시
% #2(유형 라벨)는 호환성 인자로만 받고 출력하지 않음.
\newenvironment{kfQBlock}[2]%
  {\par\smallskip\needspace{6\baselineskip}%
   \noindent\begin{minipage}{\linewidth}%
   \samepage%
   \par\noindent
   \kfQNum{#1}\hspace{4pt}%
   \begin{minipage}[t]{\dimexpr\linewidth-30pt\relax}%
  }%
  {\end{minipage}\end{minipage}\par\smallskip}

% 문제 발문
\newcommand{\kfQStem}[1]{%
  \par\noindent#1\par\vspace{2pt}%
}

% 짧은 소문항 라벨 (1), (2), (3)
\newcounter{kfSubQ}
\newcommand{\kfSubQReset}{\setcounter{kfSubQ}{0}}
\newcommand{\kfSubQ}{%
  \stepcounter{kfSubQ}%
  \par\noindent\textbf{(\arabic{kfSubQ})}~\ignorespaces%
}

% ============================================================
% activity-table.tex
% 활동: 정리표 / 내용 확인표 (마크다운 표 → tabularx 변환)
% 사용:
%   % 일반 tabular (직접 컬럼 명시)
%   \begin{kfActTable}{|l|X|X|}
%     \kfTHead{항목} & \kfTHead{설명} & \kfTHead{학생 작성} \\\hline
%     ① & ... & \kfTBlank \\\hline
%   \end{kfActTable}
%
%   % adjustbox 자동 축소 (정해진 폭으로 강제)
%   \kfActTableWrap{
%     ... (\begin{tabular}...\end{tabular} 코드) ...
%   }
% ============================================================

% 사용 시 본문 마지막 행은 \\\hline 으로 끝나야 함.
% v17.1: 흰색 mask로 Canva 배경 본문 침범 차단
\newenvironment{kfActTable}[1]%
  {\par\smallskip\needspace{4\baselineskip}%
   \renewcommand{\arraystretch}{1.3}%
   \arrayrulecolor{kfRule}%
   \begin{tcolorbox}[blanker, colback=white, left=2pt, right=2pt, top=2pt, bottom=2pt]%
   \noindent\begin{adjustbox}{max width=\linewidth}%
   \begin{tabular}{#1}%
   \hline}%
  {\end{tabular}%
   \end{adjustbox}%
   \end{tcolorbox}%
   \arrayrulecolor{black}%
   \par\smallskip}

% 헤더 행 헬퍼
\newcommand{\kfTHead}[1]{\textbf{\color{kfPrimary} #1}}
\newcommand{\kfTHeadRow}[1]{\rowcolor{kfPrimaryLight!50}\textbf{\color{kfPrimary} #1}}

% 빈 셀 (학생 작성용)
\newcommand{\kfTBlank}{\rule[-1.6em]{0pt}{3.4em}}
\newcommand{\kfTBlankSm}{\rule[-1.0em]{0pt}{2.4em}}

% adjustbox 강제 축소 래퍼 — Python에서 변환된 raw tabular 블록을 감쌀 때 사용
\newcommand{\kfActTableWrap}[1]{%
  \par\smallskip\needspace{4\baselineskip}%
  \begin{tcolorbox}[blanker, colback=white, left=2pt, right=2pt, top=2pt, bottom=2pt]%
  \noindent\begin{adjustbox}{max width=\linewidth, center}%
  #1%
  \end{adjustbox}%
  \end{tcolorbox}\par\smallskip%
}

% ============================================================
% activity-vocab.tex
% 어휘 학습 표 (3열: 번호 - 뜻 - 빈칸)
% 사용:
%   \begin{kfVocab}
%     \kfVocabItem{1}{눈으로 보고 마음으로 깨달음}
%     \kfVocabItem{2}{사물의 이치를 따져 헤아림}
%   \end{kfVocab}
% ============================================================

% adjustbox로 폭 강제 + 일반 tabular + p{} 열 사용
\newenvironment{kfVocab}%
  {\par\smallskip\needspace{4\baselineskip}%
   \renewcommand{\arraystretch}{1.35}%
   \arrayrulecolor{kfRule}%
   \noindent\begin{adjustbox}{max width=\linewidth, center}%
   \begin{tabular}{|>{\centering\arraybackslash}m{1.0cm}|m{8.5cm}|>{\centering\arraybackslash}m{4.0cm}|}%
   \hline
   \rowcolor{kfPrimaryLight!50}%
   \multicolumn{1}{|c|}{\textbf{\color{kfPrimary}번호}} &
   \multicolumn{1}{c|}{\textbf{\color{kfPrimary}뜻풀이}} &
   \multicolumn{1}{c|}{\textbf{\color{kfPrimary}어휘 (정답 작성)}} \\\hline
   \ignorespaces}%
  {\end{tabular}%
   \end{adjustbox}%
   \arrayrulecolor{black}%
   \par\smallskip}

\newcommand{\kfVocabItem}[2]{%
  \textbf{#1} & #2 & \rule[-1.0em]{0pt}{2.4em}\\\hline%
}

% ============================================================
% activity-blank.tex
% 빈칸 채우기 활동
% 사용:
%   \begin{kfBlankList}
%     \kfBlankItem 첫 번째 문장에서 (\kfBlank)은(는) ...
%     \kfBlankItem 둘째 문장에서 ...
%   \end{kfBlankList}
% ============================================================

\newenvironment{kfBlankList}%
  {\par\smallskip%
   \begin{enumerate}[leftmargin=20pt, itemsep=4pt, topsep=2pt,
                    label=\textbf{\arabic*.}, labelsep=6pt]}%
  {\end{enumerate}\par\smallskip}

\newcommand{\kfBlankItem}{\item}

% 텍스트 안에 박스형 빈칸 (단어 1개 채우기)
\newcommand{\kfBlankBox}[1][2.4cm]{%
  \tikz[baseline=-0.5ex]{%
    \draw[kfRule, line width=0.4pt] (0,0) -- ++(#1,0);%
  }%
}

% ============================================================
% activity-ox.tex (v15)
% O/X 표기 활동 — 그래픽 디자인
% 사용:
%   \begin{kfOXList}
%     \kfOXItem 글쓴이는 자연을 예찬하고 있다.\kfOX
%     \kfOXItem 1연과 4연은 수미상관 구조이다.\kfOX
%   \end{kfOXList}
%
% \kfOX = 우측에 [O] [X] 두 박스 (학생이 하나 동그라미)
% ============================================================

\newenvironment{kfOXList}%
  {\par\smallskip%
   \begin{enumerate}[leftmargin=20pt, itemsep=5pt, topsep=2pt,
                    label=\textbf{\arabic*.}, labelsep=6pt]}%
  {\end{enumerate}\par\smallskip}

\newcommand{\kfOXItem}{\item}

% v16: O/X 그래픽 박스 키움 (18×14 → 24×20pt) + 영역색 강조
\newcommand{\kfOX}{%
  \hfill\nolinebreak
  \tikz[baseline=-0.9ex]{%
    \node[draw=kfPrimary, fill=kfPrimary!5, line width=0.6pt, rounded corners=3pt,
          minimum width=24pt, minimum height=20pt, inner sep=0pt,
          font=\normalsize\bfseries\color{kfPrimary}] (ob) {O};%
    \node[draw=kfMute, fill=kfMute!5, line width=0.6pt, rounded corners=3pt,
          minimum width=24pt, minimum height=20pt, inner sep=0pt,
          font=\normalsize\bfseries\color{kfMute},
          right=4pt of ob] (xb) {X};%
  }%
}

% ============================================================
% activity-connect.tex
% 좌우 선 긋기 활동 (2단 양식 + 중앙 여백)
% 사용:
%   \begin{kfConnect}
%     \kfPair{사르트르}{실존은 본질에 앞선다}
%     \kfPair{니체}{초인 사상}
%     \kfPair{하이데거}{현존재}
%   \end{kfConnect}
% ============================================================

\newenvironment{kfConnect}%
  {\par\smallskip%
   \renewcommand{\arraystretch}{1.6}%
   \begin{tabular}{@{}>{\raggedleft\arraybackslash}m{4.5cm}
                     @{\hspace{2.5cm}}
                     >{\raggedright\arraybackslash}m{6.5cm}@{}}}%
  {\end{tabular}\par\smallskip}

\newcommand{\kfPair}[2]{%
  \tikz[baseline]{\node[draw=kfRule, fill=kfPrimaryLight!30, rounded corners=2pt,
        inner sep=4pt, font=\small]{#1};} \tikz[baseline]{\node[circle, draw=kfMute, fill=kfPaper, inner sep=1.5pt]{};}
  &
  \tikz[baseline]{\node[circle, draw=kfMute, fill=kfPaper, inner sep=1.5pt]{};} \tikz[baseline]{\node[draw=kfRule, fill=kfNoteBg, rounded corners=2pt,
        inner sep=4pt, font=\small]{#2};} \\
}

% ============================================================
% activity-writing.tex
% 글쓰기 활동 (큰 노트 영역)
% 사용:
%   \kfWriting{독후감}{12}     % 라벨 + 12줄 큰 노트
%   \kfWritingPrompt{프롬프트 텍스트}
% ============================================================

\newcommand{\kfWritingPrompt}[1]{%
  \par\smallskip%
  \begin{tcolorbox}[
    enhanced, breakable=false,
    colback=kfPrimaryLight!30, colframe=kfPrimary,
    boxrule=0.4pt, arc=3pt,
    left=\kfBoxPad, right=\kfBoxPad, top=\kfBoxPad, bottom=\kfBoxPad,
  ]%
  \small\textbf{\color{kfPrimary}글쓰기 안내}\\[2pt]%
  #1
  \end{tcolorbox}\par\smallskip%
}

\newcommand{\kfWriting}[2]{%
  \par\smallskip\needspace{\numexpr#2+2\relax\baselineskip}%
  \begin{tcolorbox}[
    enhanced, breakable=false,
    colback=kfPaper, colframe=kfRule,
    boxrule=0.4pt, arc=2pt,
    left=\kfBoxPad, right=\kfBoxPad, top=\kfBoxPad, bottom=\kfBoxPad,
    title={\footnotesize\bfseries\color{kfMute} #1},
    coltitle=kfMute, colbacktitle=kfPaper,
    attach boxed title to top left={yshift=-1.5mm, xshift=6pt},
    boxed title style={colframe=kfRule, boxrule=0.3pt, arc=1pt}
  ]%
  \foreach \i in {1,...,#2}{%
    \textcolor{kfRule}{\rule{\linewidth}{0.3pt}}\\[7pt]%
  }%
  \end{tcolorbox}\par\smallskip%
}

% ============================================================
% activity-structure.tex
% 구조도 (트리 + 빈칸)
% 사용:
%   \begin{kfStructure}
%     \kfNode{0}{서론}{문제 제기}
%     \kfNode{1}{본론}{(\,\quad\,)}
%     \kfNode{1}{본론}{근거 2}
%     \kfNode{0}{결론}{주장 강화}
%   \end{kfStructure}
% ============================================================

\newenvironment{kfStructure}%
  {\par\smallskip\needspace{6\baselineskip}%
   \begin{tcolorbox}[
     enhanced, breakable=false,
     colback=kfPaper, colframe=kfRule,
     boxrule=0.4pt, arc=2pt,
     left=\kfBoxPad, right=\kfBoxPad, top=\kfBoxPad, bottom=\kfBoxPad,
     title={\footnotesize\bfseries\color{kfMute} 글의 구조},
     coltitle=kfMute, colbacktitle=kfPrimaryLight!40,
     attach boxed title to top left={yshift=-1.5mm, xshift=6pt},
     boxed title style={colframe=kfRule, boxrule=0.3pt, arc=1pt}
   ]%
   \begin{tabbing}
     \hspace{1.4cm}\=\hspace{1.4cm}\=\hspace{8cm}\kill}%
  {\end{tabbing}\end{tcolorbox}\par\smallskip}

% 노드: #1=들여쓰기 단계 (0~3), #2=라벨, #3=내용
\newcommand{\kfNode}[3]{%
  \ifcase#1\relax%
    \>\>\textbf{\color{kfPrimary} ▶ #2}\quad #3\\
  \or
    \>\>\quad\textbf{\color{kfMute} ◦ #2}\quad #3\\
  \or
    \>\>\quad\quad\textbf{\color{kfMute} - #2}\quad #3\\
  \or
    \>\>\quad\quad\quad\textbf{\color{kfMute} \textperiodcentered\ #2}\quad #3\\
  \fi
}

% 빈칸 노드
\newcommand{\kfNodeBlank}[2]{\kfNode{#1}{#2}{(\hspace{2.5cm})}}

% ============================================================
% activity-sentence-reading.tex
% 문장 독해 활동 — 문장 박스 1개 + 그에 묶인 문제 N개를 한 minipage로
% 사용:
%   \begin{kfSentenceUnit}{1}{"바람은 내 귀에 속삭이며..."}
%     \kfSentQ{(1)}{서술어 '속삭이며'의 주어는?}
%     \begin{kfChoices}
%       \kfChoice 바람은
%       ...
%     \end{kfChoices}
%   \end{kfSentenceUnit}
% ============================================================

% kfSentenceBox 매크로는 passage-note.tex에 정의됨

\newenvironment{kfSentenceUnit}[2]%
  {\par\smallskip\needspace{6\baselineskip}%
   \noindent\begin{minipage}{\linewidth}%
   \samepage%
   \kfSentenceBox{문장 #1}{#2}%
   \par\smallskip%
  }%
  {\end{minipage}\par\smallskip}

% 같은 문장에 묶인 소문항 (문장 안내 + 발문)
\newcommand{\kfSentQ}[2]{%
  \par\noindent\textbf{#1}~#2\par\smallskip%
}

% 헤더 없이 문장만 있는 묶음 (sentence_ref가 없는 일반 문항)
\newenvironment{kfSentencelessUnit}%
  {\par\smallskip\noindent\begin{minipage}{\linewidth}\samepage}%
  {\end{minipage}\par\smallskip}

% ============================================================
% chapter-cover.tex (v6)
% 챕터 진입 페이지 (표지) — 하단에 영역 목차 추가
% 사용:
%   \kfChapterCover{1}{근대성과 자아}{Chapter 1}{이번 챕터에서는 ...}
%   \begin{kfChapterAreaList}
%     \kfChapterAreaItem{어휘}{kfAreaVocab}{핵심 어휘 12개}
%     ...
%   \end{kfChapterAreaList}
%
%   영역 목차가 없을 때는 \kfChapterCoverEnd 호출
% ============================================================

\newcommand{\kfChapterCover}[4]{%
  \clearpage%
  \thispagestyle{kfBlank}%
  \kfMarkChapter{Chapter #1. #2}%
  \kfChapterCoverBg%
  % v18.5: 챕터 표지 우상단에 로고
  \begin{tikzpicture}[remember picture, overlay]
    \node[anchor=north east, inner sep=0pt] at ($(current page.north east)+(-18mm,-18mm)$) {\kfLogoMd};
  \end{tikzpicture}%
  \vspace*{20mm}%
  \noindent
  \begin{tikzpicture}[remember picture, overlay]
    % 좌측 액센트 바
    \fill[kfPrimary] (current page.west) ++(12mm,25mm) rectangle ++(5mm, -50mm);
    % 챕터 번호 배지 (이중 원, 약간 작게)
    \node[anchor=north west, inner sep=0pt] at ($(current page.north west)+(24mm,-45mm)$) {%
      \begin{tikzpicture}
        \draw[kfPrimary, line width=1pt] (0,0) circle (10mm);
        \draw[kfPrimary, line width=0.4pt] (0,0) circle (8mm);
        \node[font=\normalsize\bfseries\color{kfPrimary}] at (0,2mm) {CHAPTER};
        \node[font=\LARGE\bfseries\color{kfPrimary}] at (0,-3mm) {#1};
      \end{tikzpicture}%
    };
  \end{tikzpicture}%
  \vspace{42mm}%
  \par\noindent\hspace{32mm}%
  \begin{minipage}{0.65\linewidth}%
    {\footnotesize\color{kfMute} #3}\\[4pt]
    {\LARGE\bfseries\color{kfPrimary} #2}\\[8pt]
    {\color{kfRule}\rule{50mm}{0.5pt}}\\[6pt]
    {\small\color{kfInk} #4}
  \end{minipage}%
  \par\vspace{14mm}%
}

% v6 / v18.5 / v18.6: 챕터 표지의 영역 목차 — 흰색 반투명 박스
% v18.6: 배경 가로선/도형과 안 겹치도록 TikZ overlay로 우하단 절대 위치
\makeatletter
% 영역 항목들을 모아 둘 박스
\newsavebox{\kf@arealistbox}
\newenvironment{kfChapterAreaList}{%
  \begin{lrbox}{\kf@arealistbox}%
  \begin{minipage}{0.62\linewidth}%
    \begin{tcolorbox}[%
      enhanced, colback=white, colframe=white,
      opacityback=0.92, opacityframe=0,
      boxrule=0pt, arc=4pt,
      width=\linewidth,
      top=8pt, bottom=8pt, left=10pt, right=10pt,
    ]%
      {\footnotesize\bfseries\color{kfMute} 이 챕터의 영역}\par\vspace{4pt}%
      {\color{kfRule}\hrule height 0.3pt}\par\vspace{4pt}%
      \begin{itemize}[leftmargin=14pt, itemsep=2pt, topsep=0pt, label={\color{kfPrimary!70}\textbullet}]%
}{%
      \end{itemize}%
    \end{tcolorbox}%
  \end{minipage}%
  \end{lrbox}%
  % v18.7: 배경 상단 가로선 ~ 하단 가로선 사이의 중간 영역에 배치
  % 가로선이 내용을 가리지 않도록 박스 중심이 두 선 사이 중점에 위치
  \begin{tikzpicture}[remember picture, overlay]%
    \node[anchor=center, inner sep=0pt]
      at ($(current page.south)+(0mm,90mm)$)
      {\usebox{\kf@arealistbox}};%
  \end{tikzpicture}%
  \vfill%
  \noindent\hfill\kfSeriesBadge\hspace{12mm}%
  \clearpage%
}
\makeatother

% 영역 항목: \kfChapterAreaItem{영역명}{kfArea색}{부제}
\makeatletter
\newcommand{\kfChapterAreaItem}[3]{%
  \item {\small\bfseries\color{#2} #1}%
  \def\kf@cci@sub{#3}%
  \ifx\kf@cci@sub\@empty\else
    {\small\color{kfMute}\, \textemdash\, #3}%
  \fi
}
\makeatother

% 영역 목록 없이 cover만 (legacy)
\newcommand{\kfChapterCoverEnd}{%
  \vfill%
  \noindent\hfill\kfSeriesBadge\hspace{12mm}%
  \clearpage%
}

% ============================================================
% area-header.tex (v16)
% 영역 시작 표제 — 시리즈별 차별화 라벨 + 큰 영역명 + 부제
% PAGE_BREAK_POLICY.md §1.4
%
% v16 (디자인 고도화 Phase 1):
%   - 시리즈별 라벨 박스 추가:
%     소쉬르 = AREA START (영역색 채움 박스)
%     프레게 = SECTION OPEN (영역색 테두리)
%     러셀   = RUSSELL (영역색 라이트 박스)
%     비트   = WITTGENSTEIN (회색 라이트 박스)
%   - 라벨 위치: 영역명 위 좌측
% ============================================================

\makeatletter

% 시리즈 라벨 텍스트
\newcommand{\kf@serieslabel}{%
  \ifcase\kf@series\relax
    AREA START%
  \or
    AREA START%
  \or
    SECTION OPEN%
  \or
    RUSSELL%
  \or
    WITTGENSTEIN%
  \fi
}

% 시리즈 라벨 박스 스타일 (#1 = 영역색)
\newcommand{\kf@labelbox}[1]{%
  \ifcase\kf@series\relax
    % 0/1 = 소쉬르: 영역색 채움
    \tikz[baseline=-2pt]{\node[fill=#1, text=white,
      rounded corners=3pt, inner xsep=6pt, inner ysep=2pt,
      font=\scriptsize\bfseries]{\kf@serieslabel};}%
  \or
    % 1 = 소쉬르: 영역색 채움
    \tikz[baseline=-2pt]{\node[fill=#1, text=white,
      rounded corners=3pt, inner xsep=6pt, inner ysep=2pt,
      font=\scriptsize\bfseries]{\kf@serieslabel};}%
  \or
    % 2 = 프레게: 영역색 테두리
    \tikz[baseline=-2pt]{\node[draw=#1, line width=0.6pt, text=#1,
      rounded corners=3pt, inner xsep=6pt, inner ysep=2pt,
      font=\scriptsize\bfseries]{\kf@serieslabel};}%
  \or
    % 3 = 러셀: 영역색 라이트 박스
    \tikz[baseline=-2pt]{\node[fill=#1!12, text=#1,
      rounded corners=2pt, inner xsep=5pt, inner ysep=1.5pt,
      font=\scriptsize\bfseries]{\kf@serieslabel};}%
  \or
    % 4 = 비트: 회색 미니 박스
    \tikz[baseline=-2pt]{\node[fill=kfMute!10, text=kfMute,
      rounded corners=2pt, inner xsep=5pt, inner ysep=1.5pt,
      font=\scriptsize\bfseries]{\kf@serieslabel};}%
  \fi
}

% v17: 시리즈+영역별 Canva 배경 자동 매핑
% \kf@hasbg=1로 set되면 \kfAreaStart에서 라벨 박스 출력 생략
\def\kf@areaVocab{어휘}\def\kf@areaGrammar{문법}\def\kf@areaConcept{개념}
\def\kf@areaLit{문학}\def\kf@areaNonlit{비문학}
\newif\ifkf@bgon
% 파일 있으면 배경 적용 + boolean true. 없으면 nothing.
\newcommand{\kfBgSet}[1]{%
  \IfFileExists{#1.pdf}{\kfPageBg{#1}\kf@bgontrue}{}%
}
\newcommand{\kf@trybg}[1]{%
  % v18.4: 영역 배경 PDF 비활성화 — 큰 도형이 본문 영역을 침범해
  % "영역 헤더 후 빈 페이지" 문제를 일으킴. 라벨 박스로만 영역 표시.
  \kf@bgonfalse%
}

% Note: 러셀(rus_*)/비트(wit_*) 배경 PDF 미존재 시 \kfPageBg가 에러 → \IfFileExists로 가드
% (현재는 파일이 모두 존재한다고 가정. 부분 제공 시 cls 수정 필요)

\newcommand{\kfAreaStart}[3]{%
  \clearpage%
  \thispagestyle{fancy}%
  \kfMarkArea{#1}%
  \kf@trybg{#1}%
  \vspace*{1mm}%
  % 배경 없으면 라텍스 라벨 박스 출력 (배경 있으면 일러스트가 라벨 역할)
  \ifkf@bgon
    \vspace*{4mm}%
  \else
    \noindent\kf@labelbox{#2}\par\vspace{4pt}%
  \fi
  % 영역명 + 부제
  \noindent
  \begin{tikzpicture}
    \node[anchor=west, inner sep=0pt] (bar) at (0,0) {\color{#2}\rule{3pt}{22pt}};
    \node[anchor=base west, inner sep=0pt, xshift=8pt, yshift=2pt] (lbl) at (bar.east |- 0,0) {%
      {\LARGE\bfseries\color{#2} #1}%
    };
    \def\kf@areasub{#3}%
    \ifx\kf@areasub\@empty\else
      \node[anchor=base west, inner sep=0pt, xshift=8pt, font=\normalsize\color{kfMute}]
        at (lbl.base east) {\,\textperiodcentered\, #3};%
    \fi
  \end{tikzpicture}%
  \par\vspace{2pt}
  {\color{#2!70}\hrule height 1pt}%
  \par\vspace{1pt}%
  {\color{kfRule}\hrule height 0.3pt}%
  \par\vspace{4pt}%
  \needspace{6\baselineskip}\nopagebreak[4]%
  \global\kf@areaJustOpenedtrue% v18.5: 첫 컨텐츠 needspace 무시
}
\makeatother

% 시리즈/영역 단축 매크로
\newcommand{\kfStartVocab}[1]{\kfAreaStart{어휘}{kfAreaVocab}{#1}}
\newcommand{\kfStartGrammar}[1]{\kfAreaStart{문법}{kfAreaGrammar}{#1}}
\newcommand{\kfStartConcept}[1]{\kfAreaStart{개념}{kfAreaConcept}{#1}}
\newcommand{\kfStartLit}[1]{\kfAreaStart{문학}{kfAreaLiterature}{#1}}
\newcommand{\kfStartNonlit}[1]{\kfAreaStart{비문학}{kfAreaNonlit}{#1}}
\newcommand{\kfStartWeekly}[1]{\kfAreaStart{실력 확인}{kfAreaWeekly}{#1}}

% ============================================================
% section-header.tex
% 섹션 시작 라벨 헤더 — [지문] / [활동] / [문제] 라벨
% 좌측 액센트 바 + 라벨 + 부제 (영역 액센트 색)
%
% 사용:
%   \kfSecHeader{kfAreaLiterature}{지문}{빼앗긴 들에도 봄은 오는가 / 이상화}
%   \kfSecHeader{kfAreaConcept}{활동}{내용 확인}
%   \kfSecHeader{kfAreaNonlit}{문제}{}
%
% 헤더 직후 페이지 분할 금지: \needspace{4\baselineskip} + \nopagebreak
% ============================================================

\providecommand{\kfSecHeader}[3]{%
  \par\needspace{10\baselineskip}\filbreak%
  \vspace{3pt}%
  \noindent
  \begin{tikzpicture}
    \node[anchor=west, inner sep=0pt] (bar) at (0,0) {\color{#1}\rule{3pt}{14pt}};
    \node[anchor=west, inner sep=0pt, xshift=6pt, font=\normalsize\bfseries\color{#1}]
       (lbl) at (bar.east) {[\,#2\,]};
    \node[anchor=west, inner sep=0pt, xshift=4pt, font=\small\color{kfMute}]
       at (lbl.east) {#3};
  \end{tikzpicture}%
  \par\vspace{1pt}%
  {\color{#1!50}\hrule height 0.4pt}%
  \par\vspace{2pt}\nopagebreak%
}

% 영역별 단축 매크로
\newcommand{\kfSecPassageVocab}[1]{\kfSecHeader{kfAreaVocab}{지문}{#1}}
\newcommand{\kfSecPassageGrammar}[1]{\kfSecHeader{kfAreaGrammar}{지문}{#1}}
\newcommand{\kfSecPassageConcept}[1]{\kfSecHeader{kfAreaConcept}{지문}{#1}}
\newcommand{\kfSecPassageLit}[1]{\kfSecHeader{kfAreaLiterature}{지문}{#1}}
\newcommand{\kfSecPassageNonlit}[1]{\kfSecHeader{kfAreaNonlit}{지문}{#1}}
\newcommand{\kfSecPassageWeekly}[1]{\kfSecHeader{kfAreaWeekly}{지문}{#1}}

\newcommand{\kfSecActVocab}[1]{\kfSecHeader{kfAreaVocab}{활동}{#1}}
\newcommand{\kfSecActGrammar}[1]{\kfSecHeader{kfAreaGrammar}{활동}{#1}}
\newcommand{\kfSecActConcept}[1]{\kfSecHeader{kfAreaConcept}{활동}{#1}}
\newcommand{\kfSecActLit}[1]{\kfSecHeader{kfAreaLiterature}{활동}{#1}}
\newcommand{\kfSecActNonlit}[1]{\kfSecHeader{kfAreaNonlit}{활동}{#1}}
\newcommand{\kfSecActWeekly}[1]{\kfSecHeader{kfAreaWeekly}{활동}{#1}}

\newcommand{\kfSecQVocab}[1]{\kfSecHeader{kfAreaVocab}{문제}{#1}}
\newcommand{\kfSecQGrammar}[1]{\kfSecHeader{kfAreaGrammar}{문제}{#1}}
\newcommand{\kfSecQConcept}[1]{\kfSecHeader{kfAreaConcept}{문제}{#1}}
\newcommand{\kfSecQLit}[1]{\kfSecHeader{kfAreaLiterature}{문제}{#1}}
\newcommand{\kfSecQNonlit}[1]{\kfSecHeader{kfAreaNonlit}{문제}{#1}}
\newcommand{\kfSecQWeekly}[1]{\kfSecHeader{kfAreaWeekly}{문제}{#1}}


\begin{document}

\thispagestyle{kfBlank}
\kfBookCoverBg
\vspace*{\kfBookCoverVspace}
\begin{center}
{\kfLogoLg}\\[14pt]
{\fontsize{72}{84}\selectfont\bfseries\kfBookCoverColor 소쉬르}\\[18pt]
{\fontsize{28}{32}\selectfont\kfBookCoverSubColor \textsc{Level} \textbf{1}}\\[22pt]
{\large\kfBookCoverSubColor 제 1 권 \quad\textbar\quad 챕터 1\textendash{} 4}
\end{center}
\vfill
\hfill\kfSeriesBadge\hspace{15mm}
\clearpage
\kfChapterCover{1}{소쉬르1 (챕터1)}{Chapter 1}{이번 챕터의 학습 내용을 시작합니다.}
\begin{kfChapterAreaList}
\kfChapterAreaItem{어휘}{kfAreaVocab}{핵심 어휘 익히기}
\kfChapterAreaItem{개념}{kfAreaConcept}{배경지식 넓히기: 우리 몸을 지켜 주는 '감각'}
\kfChapterAreaItem{문학}{kfAreaLiterature}{「눈은 반짝반짝 세상을 봐요」}
\kfChapterAreaItem{비문학}{kfAreaNonlit}{동물들의 슈퍼 감각}
\kfChapterAreaItem{실력확인}{kfAreaWeekly}{}
\end{kfChapterAreaList}

\kfStartVocab{핵심 어휘 익히기}


\begin{kfVocab}
\kfVocabItem{1}{눈으로 보는 감각}
\kfVocabItem{2}{귀로 듣는 감각}
\kfVocabItem{3}{코로 냄새 맡는 감각}
\kfVocabItem{4}{혀로 맛보는 감각}
\kfVocabItem{5}{피부로 느끼는 감각}
\kfVocabItem{6}{몸으로 느끼는 것}
\end{kfVocab}


\kfStartConcept{배경지식 넓히기: 우리 몸을 지켜 주는 '감각'}


\kfPassageNoteNT{%
우리가 느끼는 감각은 세상을 즐겁게 탐험하도록 도와주기도 하지만, 위험한 상황에서 우리 몸을 지켜 주는 '보디가드' 역할도 해요.

코(후각)의 경고: "으악, 쉰 냄새가 나!" 상한 음식 냄새를 맡게 해서, 배가 아프지 않게 먹지 말라고 알려 줘요.

피부(촉각)의 경고: "앗, 뜨거워!" 뜨거운 주전자에 손이 닿으면 깜짝 놀라 손을 떼게 해서, 화상을 입지 않게 도와줘요.

귀(청각)의 경고: "빵빵! 비키세요!" 뒤에서 오는 자동차 소리를 듣고, 다치지 않게 안전한 곳으로 피하게 해 줘요.
\par\medskip\begin{itemize}[leftmargin=18pt, itemsep=2pt, topsep=2pt, label=\textbullet]
\item 코(후각): 상한 음식 냄새를 맡고 멀리합니다.
\item 피부(촉각): 뜨거운 주전자에서 손을 뗍니다.
\item 귀(청각): 자동차 경적을 듣고 안전한 곳으로 피합니다.
\end{itemize}
}

\kfActivityBg \par\kfMaybeNeedspace{24\baselineskip}\noindent{\kfTipMark\,\,}{\small\color{kfInk} 다음 단어와 알맞은 설명을 연결해 보세요.}\par\nopagebreak[4]\smallskip\nopagebreak[4]
\begin{kfConnect}
\kfPair{눈}{사물을 봐요}
\kfPair{코}{소리를 들어요}
\kfPair{입}{만져서 느껴요}
\kfPair{귀}{냄새를 맡아요}
\kfPair{손}{맛을 느껴요}
\end{kfConnect}

\kfActivityBg \par\kfMaybeNeedspace{24\baselineskip}\noindent{\kfTipMark\,\,}{\small\color{kfInk} 다음 문장이 옳으면 O, 옳지 않으면 X를 표시하세요.}\par\nopagebreak[4]\smallskip\nopagebreak[4]
\begin{kfOXList}
\kfOXItem 감각은 우리를 즐겁게 해 줄 뿐만 아니라, 위험할 때 몸을 지켜 주기도 해요. \kfFillBlank[12mm]{} \kfOX
\kfOXItem 뜨거운 주전자에 손이 닿았을 때 위험을 알려 주는 감각은 '후각'이에요. \kfFillBlank[12mm]{} \kfOX
\kfOXItem 자동차 소리를 듣고 안전한 곳으로 피하게 해 주는 것은 '청각' 덕분이에요. \kfFillBlank[12mm]{} \kfOX
\end{kfOXList}

\kfActivityBg \par\kfMaybeNeedspace{24\baselineskip}\noindent{\kfTipMark\,\,}{\small\color{kfInk} 다음 문장에 어울리는 감각을 나타내는 말을 고르세요.}\par\nopagebreak[4]\smallskip\nopagebreak[4]
\begin{kfBlankList}
\kfBlankItem 눈으로 별을 보니 ( 반짝반짝 / 킁킁 / 쫑긋쫑긋 ) 빛났어요.
\kfBlankItem 코로 꽃 냄새를 ( 통통통 / 킁킁 / 냠냠 ) 맡았어요.
\kfBlankItem 귀를 ( 살살살 / 쫑긋쫑긋 / 반짝반짝 ) 세우고 소리를 들었어요.
\end{kfBlankList}


\kfStartLit{「눈은 반짝반짝 세상을 봐요」}


\kfPassageNote{}{%
눈은 반짝반짝 세상을 봐요 \\[4pt]
코는 킁킁킁 냄새를 맡아요 \\[4pt]
귀는 쫑긋쫑긋 소리를 들어요 \\[4pt]
입은 냠냠냠 맛을 느껴요 \\[4pt]
손은 살살살 만져 봐요 \\[4pt]
발은 통통통 걸어가요 \\[14pt]
우리 몸은 정말 신기해요 \\[4pt]
내 몸은 소중해요
}


\kfActivityBg \par\kfMaybeNeedspace{18\baselineskip}\noindent{\kfTipMark\,\,}{\small\color{kfInk} 주어진 안내에 맞춰 자유롭게 글로 표현해 보세요.}\par\nopagebreak[4]\smallskip\nopagebreak[4]
\kfWritingPrompt{동요에 나오지 않은 다른 신체 부위를 넣거나 재미있는 흉내 내는 말(의성어·의태어)을 넣어 나만의 노랫말을 만들어 보세요.}
\kfWriting{답안 작성}{8}

\kfSecHeader{kfAreaLiterature}{문제}{}

\begin{kfQBlock}{01}{객관식}
\kfQStem{이 동요에서 '반짝반짝' 세상을 보는 신체 부위는 무엇인가요?}
\begin{kfChoices}
\kfChoice 코
\kfChoice 귀
\kfChoice 눈
\kfChoice 입
\end{kfChoices}
\end{kfQBlock}
\begin{kfQBlock}{02}{객관식}
\kfQStem{동요에서 발이 하는 행동을 나타낸 말은 무엇인가요?}
\begin{kfChoices}
\kfChoice 살살살
\kfChoice 통통통
\kfChoice 킁킁킁
\kfChoice 쫑긋쫑긋
\end{kfChoices}
\end{kfQBlock}
\begin{kfQBlock}{03}{객관식}
\kfQStem{이 동요에서 '냠냠냠 맛을 느낀다'고 한 신체 부위는 무엇인가요?}
\begin{kfChoices}
\kfChoice 코
\kfChoice 입
\kfChoice 손
\kfChoice 귀
\end{kfChoices}
\end{kfQBlock}
\begin{kfQBlock}{04}{서술형}
\kfQStem{이 동요를 읽고 나서, 여러분이 가장 고마운 신체 부위는 무엇인가요? 그 이유도 함께 써 보세요.}
\kfAnswerNote{답안}{4}
\end{kfQBlock}
\par\noindent\textit{\small * 나는 \kfFillBlank[12mm]{} 이(가) 가장 고마워요. * 왜냐하면 \_\_\_\_\_\_\_\_\_\_\_\_\_\_\_\_\_\_\_\_\_\_\_\_\_\_\_\_\_\_\_\_\_\_\_\_\_\_\_\_}


\kfStartNonlit{동물들의 슈퍼 감각}


\kfPassageNote{동물들의 슈퍼 감각}{%
우리 몸에는 세상을 보고, 냄새 맡고, 들을 수 있게 해 주는 고마운 눈, 코, 귀가 있어요. 그런데 동물 친구들 중에는 사람보다 훨씬 더 뛰어난 ‘슈퍼 감각’을 가진 친구들이 있답니다.

먼저 매는 아주 특별한 눈을 가지고 있어요. 매는 높은 하늘을 날아다니지만, 땅 위에 있는 아주 작은 토끼도 금방 찾아낼 수 있습니다. 사람보다 무려 여덟 배나 더 멀리, 더 자세히 볼 수 있기 때문이에요.

개는 훌륭한 코를 가지고 있어요. 킁킁! 개는 냄새를 아주 잘 맡아서 잃어버린 물건을 찾거나, 눈에 보이지 않는 범인을 냄새로 찾아내기도 합니다. 경찰관 아저씨를 돕는 경찰견은 바로 이 뛰어난 냄새 맡는 능력 덕분에 멋진 일을 할 수 있는 거예요.

마지막으로 토끼는 길쭉한 귀를 가지고 있어요. 토끼의 귀가 긴 이유는 무엇일까요? 바로 아주 작은 소리도 놓치지 않고 잘 듣기 위해서예요. 토끼는 귀를 이리저리 움직여서 사방에서 들려오는 소리를 모을 수 있답니다. 동물들의 능력은 정말 신기하지 않나요?
}


\kfActivityBg \par\needspace{15\baselineskip}
\kfSecHeader{kfAreaNonlit}{활동}{문장 독해}
\par\kfMaybeNeedspace{32\baselineskip}\noindent{\kfTipMark\,\,}{\small\color{kfInk} 각 문장을 읽고 물음에 답하세요.}\par\nopagebreak[4]\smallskip\nopagebreak[4]
\begin{kfSentenceUnit}{1}{눈은 빛을 받아들여요}
\kfSentQ{1.}{(1) 이 문장에서 '무엇이'에 해당하는 것은 무엇인가요?}
\begin{kfChoices}
\kfChoice 눈은
\kfChoice 빛을
\kfChoice 받아들여요
\end{kfChoices}
\kfSentQ{1.}{(2) 이 문장에서 '어찌하다'에 해당하는 것은 무엇인가요?}
\begin{kfChoices}
\kfChoice 눈은
\kfChoice 빛을
\kfChoice 받아들여요
\end{kfChoices}
\end{kfSentenceUnit}

\begin{kfSentenceUnit}{2}{코는 먼지를 깨끗하게}
\kfSentQ{2.}{(1) 이 문장에서 '무엇이'에 해당하는 것은 무엇인가요?}
\begin{kfChoices}
\kfChoice 코는
\kfChoice 먼지를
\kfChoice 깨끗하게
\end{kfChoices}
\kfSentQ{2.}{(2) 이 문장에서 '무엇을'에 해당하는 것은 무엇인가요?}
\begin{kfChoices}
\kfChoice 코는
\kfChoice 먼지를
\kfChoice 깨끗하게
\end{kfChoices}
\end{kfSentenceUnit}

\kfActivityBg \par\needspace{15\baselineskip}
\kfSecHeader{kfAreaNonlit}{활동}{문장 독해}
\par\kfMaybeNeedspace{32\baselineskip}\noindent{\kfTipMark\,\,}{\small\color{kfInk} 각 문장을 읽고 물음에 답하세요.}\par\nopagebreak[4]\smallskip\nopagebreak[4]
\begin{kfSentenceUnit}{1}{눈은 빛을 받아들여요}
\kfSentQ{1.}{(1) 이 문장에서 '무엇이'에 해당하는 것은 무엇인가요?}
\begin{kfChoices}
\kfChoice 눈은
\kfChoice 빛을
\kfChoice 받아들여요
\end{kfChoices}
\kfSentQ{1.}{(2) 이 문장에서 '어찌하다'에 해당하는 것은 무엇인가요?}
\begin{kfChoices}
\kfChoice 눈은
\kfChoice 빛을
\kfChoice 받아들여요
\end{kfChoices}
\end{kfSentenceUnit}

\begin{kfSentenceUnit}{2}{코는 먼지를 깨끗하게}
\kfSentQ{2.}{(1) 이 문장에서 '무엇이'에 해당하는 것은 무엇인가요?}
\begin{kfChoices}
\kfChoice 코는
\kfChoice 먼지를
\kfChoice 깨끗하게
\end{kfChoices}
\kfSentQ{2.}{(2) 이 문장에서 '무엇을'에 해당하는 것은 무엇인가요?}
\begin{kfChoices}
\kfChoice 코는
\kfChoice 먼지를
\kfChoice 깨끗하게
\end{kfChoices}
\end{kfSentenceUnit}

\kfActivityBg \par\kfMaybeNeedspace{24\baselineskip}\noindent{\kfTipMark\,\,}{\small\color{kfInk} 빈칸에 알맞은 말을 채워 보세요.}\par\nopagebreak[4]\smallskip\nopagebreak[4]
\begin{kfBlankList}
\kfBlankItem 빈칸에 알맞은 말을 <보기>에서 찾아 써 보세요.

* 매는 사람보다 8배나 더 잘 보는 \kfFillBlank[42mm]{}을 가지고 있어요. * 개는 \kfFillBlank[42mm]{}를 아주 잘 맡아서 잃어버린 물건을 찾아요. * 토끼는 긴 귀로 아주 작은 \kfFillBlank[42mm]{}도 잘 들을 수 있어요.
\begin{kfEvidence}냄새, 눈, 소리, 코, 맛\end{kfEvidence}
\end{kfBlankList}

\kfActivityBg \par\kfMaybeNeedspace{18\baselineskip}\noindent{\kfTipMark\,\,}{\small\color{kfInk} 주어진 안내에 맞춰 자유롭게 글로 표현해 보세요.}\par\nopagebreak[4]\smallskip\nopagebreak[4]
\kfWritingPrompt{만약 여러분이 동물들의 슈퍼 감각 중 하나를 가질 수 있다면, 무엇을 갖고 싶은가요? 그리고 그 감각으로 무엇을 하고 싶은지 상상해서 써 보세요.}
\kfWriting{답안 작성}{8}

\kfSecHeader{kfAreaNonlit}{문제}{}

\begin{kfQBlock}{01}{객관식}
\kfQStem{이 글에서 설명하지 \uline{\underline{않은}} 동물은 누구인가요?}
\begin{kfChoices}
\kfChoice 매
\kfChoice 개
\kfChoice 고양이
\kfChoice 토끼
\end{kfChoices}
\end{kfQBlock}
\begin{kfQBlock}{02}{객관식}
\kfQStem{토끼의 귀가 길쭉한 이유는 무엇인가요?}
\begin{kfChoices}
\kfChoice 빠르게 하늘을 날기 위해서
\kfChoice 음식 냄새를 잘 맡기 위해서
\kfChoice 작은 소리도 놓치지 않으려고
\kfChoice 사방을 멀리까지 보기 위해서
\end{kfChoices}
\end{kfQBlock}
\begin{kfQBlock}{03}{객관식}
\kfQStem{매는 사람보다 눈이 얼마나 더 좋다고 했나요?}
\begin{kfChoices}
\kfChoice 2배
\kfChoice 4배
\kfChoice 8배
\kfChoice 10배
\end{kfChoices}
\end{kfQBlock}


\kfStartWeekly{실력확인 영역 학습}


\noindent\textit{\small 제한시간: 20분 / 총 10문항}\par\medskip
\par\noindent\textbf{[3\textasciitilde{}5] 다음 글을 읽고 물음에 답하시오.}\par\medskip
\kfPassageNote{지문 3}{%
내 몸은 소중해 \\[14pt]
눈은 ㉠\uline{반짝반짝} 세상을 봐요 \\[4pt]
코는 킁킁킁 냄새를 맡아요 \\[4pt]
귀는 쫑긋쫑긋 소리를 들어요 \\[4pt]
입은 냠냠냠 맛을 느껴요 \\[4pt]
손은 살살살 만져 봐요 \\[4pt]
발은 통통통 걸어가요 \\[14pt]
우리 몸은 정말 신기해요 \\[4pt]
내 몸은 소중해요
}
\par\noindent\textbf{[6\textasciitilde{}8] 다음 글을 읽고 물음에 답하시오.}\par\medskip
\kfPassageNote{지문 6}{%
동물들의 슈퍼 감각

우리 몸에는 세상을 보고, 냄새 맡고, 들을 수 있게 해 주는 고마운 눈, 코, 귀가 있어요. 그런데 동물 친구들 중에는 사람보다 훨씬 더 뛰어난 ‘슈퍼 감각’을 가진 친구들이 있답니다.

먼저 ( ⓐ )은/는 아주 특별한 눈을 가지고 있어요. 높은 하늘을 날아다니지만, 땅 위에 있는 아주 작은 토끼도 금방 찾아낼 수 있습니다. 사람보다 무려 여덟 배나 더 멀리, 더 자세히 볼 수 있기 때문이에요.

개는 훌륭한 코를 가지고 있어요. 킁킁! 개는 냄새를 아주 잘 맡아서 잃어버린 물건을 찾거나, 눈에 보이지 않는 범인을 냄새로 찾아내기도 합니다.

마지막으로 토끼는 길쭉한 귀를 가지고 있어요. 토끼의 귀가 긴 이유는 바로 아주 작은 소리도 놓치지 않고 잘 듣기 위해서예요. 토끼는 귀를 이리저리 움직여서 사방에서 들려오는 소리를 모을 수 있답니다.
}
\begin{kfQBlock}{01}{객관식}
\kfQStem{다음 중 맛을 느끼는 감각 기관은 무엇인가요?}
\begin{kfChoices}
\kfChoice 눈
\kfChoice 귀
\kfChoice 입
\kfChoice 코
\end{kfChoices}
\end{kfQBlock}

\begin{kfQBlock}{02}{객관식}
\kfQStem{다음 설명에 알맞은 감각의 이름은 무엇인가요?}
\begin{kfEvidence}"코로 냄새를 맡는 감각이에요. 상한 음식 냄새를 맡아 배탈이 나지 않게 도와주기도 해요."\end{kfEvidence}
\begin{kfChoices}
\kfChoice 시각
\kfChoice 청각
\kfChoice 후각
\kfChoice 촉각
\end{kfChoices}
\end{kfQBlock}

\begin{kfQBlock}{03}{객관식}
\kfQStem{㉠과 같이 모양이나 소리를 흉내 내는 말이 \uline{\underline{아닌}} 것은 무엇인가요?}
\begin{kfChoices}
\kfChoice 킁킁킁
\kfChoice 냠냠냠
\kfChoice 소중해
\kfChoice 통통통
\end{kfChoices}
\end{kfQBlock}

\begin{kfQBlock}{04}{객관식}
\kfQStem{다음 중 이 동시에 나오지 \underline{않은} 신체 부위는 어디인가요?}
\begin{kfChoices}
\kfChoice 귀
\kfChoice 발
\kfChoice 입
\kfChoice 머리
\end{kfChoices}
\end{kfQBlock}

\begin{kfQBlock}{05}{객관식}
\kfQStem{빈칸 ⓐ에 들어갈 알맞은 동물은 누구인가요?}
\begin{kfChoices}
\kfChoice 매
\kfChoice 호랑이
\kfChoice 강아지
\kfChoice 고양이
\end{kfChoices}
\end{kfQBlock}

\begin{kfQBlock}{06}{객관식}
\kfQStem{개가 잃어버린 물건을 잘 찾을 수 있는 이유는 무엇인가요?}
\begin{kfChoices}
\kfChoice 눈이 좋아서 멀리까지 볼 수 있기 때문에
\kfChoice 귀가 커서 소리를 잘 듣기 때문에
\kfChoice 코로 냄새를 아주 잘 맡기 때문에
\kfChoice 발이 빨라서 금방 달려갈 수 있기 때문에
\end{kfChoices}
\end{kfQBlock}

\begin{kfQBlock}{07}{객관식}
\kfQStem{다음 중 '입'이 하는 일이 \uline{\underline{아닌}} 것은 무엇인가요?}
\begin{kfChoices}
\kfChoice 음식을 씹어서 먹는다.
\kfChoice 맛을 느낀다.
\kfChoice 친구와 대화할 때 말을 한다.
\kfChoice 반짝반짝 빛을 받아들인다.
\end{kfChoices}
\end{kfQBlock}

\begin{kfQBlock}{08}{서술형}
\kfQStem{이 동시에서 '냠냠냠' 맛을 느끼는 신체 부위는 무엇인지 찾아 쓰세요.}
\kfAnswerNote{답안}{4}
\end{kfQBlock}

\begin{kfQBlock}{09}{서술형}
\kfQStem{토끼의 귀가 길쭉한 이유는 무엇인가요? 글의 내용을 바탕으로 한 문장으로 쓰세요.}
\kfAnswerNote{답안}{4}
\end{kfQBlock}

\begin{kfQBlock}{10}{서술형}
\kfQStem{다음 문장에서 '무엇을'에 해당하는 말을 찾아 쓰세요.}
\begin{kfEvidence}"코는 먼지를 깨끗하게 걸러 줘요."\end{kfEvidence}
\kfAnswerNote{답안}{4}
\end{kfQBlock}




\clearpage

\kfChapterCover{2}{소쉬르1 (챕터2)}{Chapter 2}{이번 챕터의 학습 내용을 시작합니다.}
\begin{kfChapterAreaList}
\kfChapterAreaItem{어휘}{kfAreaVocab}{핵심 어휘 익히기}
\kfChapterAreaItem{개념}{kfAreaConcept}{배경지식 넓히기: 마음의 날씨, 감정}
\kfChapterAreaItem{문학}{kfAreaLiterature}{「하하 호호 기쁠 때는」}
\kfChapterAreaItem{비문학}{kfAreaNonlit}{감정의 종류와 표현법}
\kfChapterAreaItem{실력확인}{kfAreaWeekly}{}
\end{kfChapterAreaList}

\kfStartVocab{핵심 어휘 익히기}


\begin{kfVocab}
\kfVocabItem{1}{어려운 일을 해냈을 때 마음이 꽉 찬 듯 자랑스러운 느낌}
\kfVocabItem{2}{일이 내 생각대로 되지 않아 마음이 불편하고 안타까움}
\kfVocabItem{3}{좋은 일이 생길 것 같아 가슴이 두근거리는 마음}
\kfVocabItem{4}{중요안 일을 앞두고 마음이 조마조마하고 팽팽함}
\kfVocabItem{5}{기대했던 것과 달라서 섭섭하고 아쉬운 마음}
\kfVocabItem{6}{친구의 마음을 나도 똑같이 느끼고 이해하는 것}
\end{kfVocab}


\kfStartConcept{배경지식 넓히기: 마음의 날씨, 감정}


\kfPassageNoteNT{%
우리의 마음은 마치 날씨와 같아요. 맑은 날, 비 오는 날, 천둥이 치는 날이 있는 것처럼 우리의 기분도 매일매일 변한답니다. 이것을 '감정'이라고 해요.

좋은 일이 생기면 햇살처럼 환한 '기쁨'을 느끼고, 친구와 싸우거나 물건을 잃어버리면 먹구름처럼 어두운 '슬픔'이나 '화'를 느껴요. 깜깜한 밤에는 '두려움'을 느끼기도 하죠.

모든 감정은 소중해요. 기쁜 마음뿐만 아니라 슬프거나 화나는 마음도 내가 느끼는 자연스러운 감정이니까요. 내 마음의 날씨를 잘 살피는 어린이가 되어 볼까요?
\par\medskip\begin{itemize}[leftmargin=18pt, itemsep=2pt, topsep=2pt, label=\textbullet]
\item 기쁨 — 햇살처럼 환한 마음 (예: 친구와 같이 놀 때)
\item 슬픔 — 먹구름처럼 어두운 마음 (예: 아끼던 물건을 잃어버렸을 때)
\item 두려움 — 깜깜한 밤에 느끼는 마음
\end{itemize}
}

\kfActivityBg \par\kfMaybeNeedspace{24\baselineskip}\noindent{\kfTipMark\,\,}{\small\color{kfInk} 다음 단어와 알맞은 설명을 연결해 보세요.}\par\nopagebreak[4]\smallskip\nopagebreak[4]
\begin{kfConnect}
\kfPair{뿌듯함}{친구의 마음을 나도 똑같이 느껴요}
\kfPair{설렘}{기대와 달라서 섭섭하고 아쉬워요}
\kfPair{긴장}{어려운 일을 해내서 자랑스러워요}
\kfPair{서운함}{좋은 일이 생길 것 같아 두근거려요}
\kfPair{공감}{중요한 일을 앞두고 조마조마해요}
\end{kfConnect}

\kfActivityBg \par\kfMaybeNeedspace{24\baselineskip}\noindent{\kfTipMark\,\,}{\small\color{kfInk} 다음 문장이 옳으면 O, 옳지 않으면 X를 표시하세요.}\par\nopagebreak[4]\smallskip\nopagebreak[4]
\begin{kfOXList}
\kfOXItem 사람의 마음은 날씨처럼 매일 변할 수 있어요. \kfFillBlank[12mm]{} \kfOX
\kfOXItem 화가 나거나 슬픈 마음은 나쁜 것이니 무조건 숨겨야 해요. \kfFillBlank[12mm]{} \kfOX
\kfOXItem '두려움'은 무섭거나 겁이 나는 마음을 말해요. \kfFillBlank[12mm]{} \kfOX
\end{kfOXList}

\kfActivityBg \par\kfMaybeNeedspace{24\baselineskip}\noindent{\kfTipMark\,\,}{\small\color{kfInk} 다음 상황에 어울리는 감정 단어를 골라 동그라미 하세요.}\par\nopagebreak[4]\smallskip\nopagebreak[4]
\begin{kfBlankList}
\kfBlankItem 생일 선물로 갖고 싶었던 로봇을 받았어요. ( 기쁨 / 화 )
\kfBlankItem 캄캄한 밤에 혼자 화장실에 가야 해요. ( 즐거움 / 두려움 )
\kfBlankItem 열심히 만든 블록을 친구가 무너뜨렸어요. ( 기쁨 / 화남 )
\kfBlankItem 어려운 퍼즐을 끝까지 혼자 힘으로 다 맞췄어요. ( 뿌듯함 / 속상함 )
\kfBlankItem 내일은 기다리던 소풍날이라 가슴이 콩닥거려요. ( 설렘 / 무서움 )
\kfBlankItem 놀이공원에 가기로 했는데 굵은 비가 쏟아져서 못 가게 되었어요. ( 즐거움 / 속상함 )
\kfBlankItem 피아노 대회에서 수많은 관객들 앞에 서기 바로 전이에요. ( 편안함 / 긴장됨 )
\kfBlankItem 기대하고 있었는데 엄마가 오늘 놀아 주기로 한 약속을 지키지 못했어요. ( 공감 / 서운함 )
\end{kfBlankList}


\kfStartLit{「하하 호호 기쁠 때는」}


\kfPassageNote{}{%
하하 호호 기쁠 때는 \\[4pt]
마음에 쨍하고 해가 떠요 \\[14pt]
흑흑 훌쩍 슬플 때는 \\[4pt]
마음에 주륵주륵 비가 내려요 \\[14pt]
씩씩 쿵쿵 화날 때는 \\[4pt]
마음에 우르릉 천둥이 쳐요 \\[14pt]
덜덜 깜짝 무서울 때는 \\[4pt]
마음에 캄캄한 밤이 찾아와요 \\[14pt]
내 마음은 알록달록 \\[4pt]
하루에도 몇 번씩 날씨가 바뀌어요
}


\kfActivityBg \par\kfMaybeNeedspace{18\baselineskip}\noindent{\kfTipMark\,\,}{\small\color{kfInk} 주어진 안내에 맞춰 자유롭게 글로 표현해 보세요.}\par\nopagebreak[4]\smallskip\nopagebreak[4]
\kfWritingPrompt{동시에 나온 날씨 대신 다른 색깔이나 사물로 감정을 표현해 나만의 노랫말을 만들어 보세요.}
\kfWriting{답안 작성}{8}

\kfSecHeader{kfAreaLiterature}{문제}{}

\begin{kfQBlock}{01}{객관식}
\kfQStem{이 동시에서 기쁠 때 마음에 무엇이 뜬다고 했나요?}
\begin{kfChoices}
\kfChoice 달
\kfChoice 해
\kfChoice 구름
\kfChoice 무지개
\end{kfChoices}
\end{kfQBlock}
\begin{kfQBlock}{02}{객관식}
\kfQStem{동시에서 '화날 때'의 마음을 빗대어 표현한 날씨는 무엇인가요?}
\begin{kfChoices}
\kfChoice 비
\kfChoice 눈
\kfChoice 바람
\kfChoice 천둥
\end{kfChoices}
\end{kfQBlock}
\begin{kfQBlock}{03}{객관식}
\kfQStem{다음 중 이 동시에 나오지 \underline{않은} 감정은 무엇인가요?}
\begin{kfChoices}
\kfChoice 기쁨
\kfChoice 심심함
\kfChoice 슬픔
\kfChoice 무서움
\end{kfChoices}
\end{kfQBlock}
\begin{kfQBlock}{04}{서술형}
\kfQStem{여러분은 언제 마음속에 '해'가 뜬 것처럼 기쁜가요? 경험을 떠올려 써 보세요.}
\kfAnswerNote{답안}{4}
\end{kfQBlock}


\kfStartNonlit{감정의 종류와 표현법}


\kfPassageNote{감정의 종류와 표현법}{%
사람은 누구나 다양한 감정을 느끼며 살아갑니다. 감정에는 기쁨, 즐거움, 행복함 같은 긍정적인 감정도 있고, 슬픔, 화, 짜증, 부끄러움 같은 조금 힘든 감정도 있습니다. 이 모든 감정은 우리가 살아가는데 꼭 필요한 마음의 신호입니다.

기쁘고 즐거울 때는 웃음을 짓거나 "정말 신난다!"라고 말하며 기분을 나눌 수 있어요. 반대로 슬플 때는 억지로 참지 말고 엉엉 울어도 괜찮아요. 눈물은 슬픈 마음을 씻어 내는 데 도움을 주거든요. 화가 날 때는 심호흡을 크게 '후\textasciitilde{}' 하고 내쉬거나 숫자를 세며 마음을 가라앉히는 것이 좋아요.

가장 중요한 것은 자신의 감정을 솔직하게 말하는 것입니다. 친구가 장난감을 뺏어가서 화가 났다면, 소리를 지르거나 때리는 대신 "네가 내 장난감을 가져가서 화가 났어. 돌려줘."라고 또박또박 말해야 합니다. 내 마음을 잘 표현하면 친구와 더 잘 지낼 수 있습니다.
}


\kfActivityBg \par\needspace{15\baselineskip}
\kfSecHeader{kfAreaNonlit}{활동}{문장 독해}
\par\kfMaybeNeedspace{32\baselineskip}\noindent{\kfTipMark\,\,}{\small\color{kfInk} 각 문장을 읽고 물음에 답하세요.}\par\nopagebreak[4]\smallskip\nopagebreak[4]
\begin{kfSentenceUnit}{1}{다양한 감정을 느껴요 사람은}
\kfSentQ{1.}{(1) 이 문장에서 '무엇을'에 해당하는 것은 무엇인가요?}
\begin{kfChoices}
\kfChoice 다양한 감정을
\kfChoice 느껴요
\kfChoice 사람은
\end{kfChoices}
\kfSentQ{1.}{(2) 이 문장에서 '어찌하다'에 해당하는 것은 무엇인가요?}
\begin{kfChoices}
\kfChoice 사람은
\kfChoice 다양한 감정을
\kfChoice 느껴요
\end{kfChoices}
\end{kfSentenceUnit}

\begin{kfSentenceUnit}{2}{뺏어갔어요 내 장난감을 친구가}
\kfSentQ{2.}{(1) 이 문장에서 '누가'에 해당하는 것은 무엇인가요?}
\begin{kfChoices}
\kfChoice 뺏어갔어요
\kfChoice 내 장난감을
\kfChoice 친구가
\end{kfChoices}
\kfSentQ{2.}{(2) 이 문장에서 '무엇을'에 해당하는 것은 무엇인가요?}
\begin{kfChoices}
\kfChoice 내 장난감을
\kfChoice 뺏어갔어요
\kfChoice 친구가
\end{kfChoices}
\kfSentQ{2.}{(3) 이 문장에서 '어찌하다'에 해당하는 것은 무엇인가요?}
\begin{kfChoices}
\kfChoice 친구가
\kfChoice 뺏어갔어요
\kfChoice 내 장난감을
\end{kfChoices}
\end{kfSentenceUnit}

\kfActivityBg \par\kfMaybeNeedspace{24\baselineskip}\noindent{\kfTipMark\,\,}{\small\color{kfInk} 빈칸에 알맞은 말을 채워 보세요.}\par\nopagebreak[4]\smallskip\nopagebreak[4]
\begin{kfBlankList}
\kfBlankItem 빈칸에 알맞은 말을 <보기>에서 찾아 써 보세요.

* 기쁠 때는 \kfFillBlank[42mm]{}을/를 지으며 기분을 나눠요. * 슬플 때는 억지로 참지 말고 \kfFillBlank[42mm]{}을/를 흘려도 괜찮아요. * 화가 날 때는 소리를 지르는 대신 \kfFillBlank[42mm]{}으로 표현해요.
\begin{kfEvidence}말, 눈물, 웃음, 짜증, 부끄러움\end{kfEvidence}
\end{kfBlankList}

\kfActivityBg \par\kfMaybeNeedspace{18\baselineskip}\noindent{\kfTipMark\,\,}{\small\color{kfInk} 주어진 안내에 맞춰 자유롭게 글로 표현해 보세요.}\par\nopagebreak[4]\smallskip\nopagebreak[4]
\kfWritingPrompt{친구가 실수로 내 공책을 찢었어요. 너무 속상하고 화가 나요. 이때 친구에게 어떻게 말하면 좋을까요?}
\kfWriting{답안 작성}{8}

\kfSecHeader{kfAreaNonlit}{문제}{}

\begin{kfQBlock}{01}{객관식}
\kfQStem{이 글에서 화가 날 때 마음을 가라앉히는 방법으로 알려 준 것은 무엇인가요?}
\begin{kfChoices}
\kfChoice 입을 꾹 닫고 아무 말도 하지 않는다.
\kfChoice 큰 소리로 "정말 신난다!"라고 외친다.
\kfChoice 눈물이 멈출 때까지 한참 동안 엉엉 운다.
\kfChoice 심호흡을 크게 '후\textasciitilde{}' 내쉬거나 숫자를 센다.
\end{kfChoices}
\end{kfQBlock}
\begin{kfQBlock}{02}{객관식}
\kfQStem{이 글에서 눈물에 대해 설명한 내용으로 알맞은 것은 무엇인가요?}
\begin{kfChoices}
\kfChoice 화가 났을 때만 흘리는 신호이다.
\kfChoice 슬픈 마음을 씻어 내는 데 도움을 준다.
\kfChoice 기쁘고 즐거울 때 짓는 표정이다.
\kfChoice 살아가는 데 필요 없는 마음의 신호이다.
\end{kfChoices}
\end{kfQBlock}
\begin{kfQBlock}{03}{객관식}
\kfQStem{친구가 장난감을 뺏어 가서 화가 났을 때, 이 글에서 권한 행동은 무엇인가요?}
\begin{kfChoices}
\kfChoice 웃음을 지으며 "정말 신난다!"라고 큰 소리로 외친 뒤 함께 즐거워한다.
\kfChoice "네가 장난감을 가져가서 화가 났어. 돌려줘."라고 말한다.
\kfChoice 엉엉 울어서 마음을 씻어 낸 뒤 그냥 아무 말 없이 모른 척 넘어간다.
\kfChoice 심호흡을 크게 한 뒤 숫자만 세고 친구에게 아무 말도 하지 않은 채 돌아선다.
\end{kfChoices}
\end{kfQBlock}


\kfStartWeekly{실력확인 영역 학습}


\noindent\textit{\small 제한시간: 20분 / 총 10문항}\par\medskip
\par\noindent\textbf{[3\textasciitilde{}5] 다음 글을 읽고 물음에 답하시오.}\par\medskip
\kfPassageNote{지문 3}{%
내 마음의 날씨 \\[14pt]
하하 호호 기쁠 때는 \\[4pt]
마음에 쨍하고 해가 떠요 \\[14pt]
흑흑 훌쩍 슬플 때는 \\[4pt]
마음에 주륵주륵 비가 내려요 \\[14pt]
씩씩 쿵쿵 화날 때는 \\[4pt]
마음에 우르릉 천둥이 쳐요 \\[14pt]
덜덜 깜짝 무서울 때는 \\[4pt]
마음에 캄캄한 밤이 찾아와요
}
\par\noindent\textbf{[6\textasciitilde{}8] 다음 글을 읽고 물음에 답하시오.}\par\medskip
\kfPassageNote{지문 6}{%
감정의 종류와 표현법

사람은 누구나 다양한 감정을 느끼며 살아갑니다. 기쁨 같은 긍정적인 감정도 있고, 슬픔이나 화 같은 힘든 감정도 있습니다. 이 모든 감정은 우리가 살아가는데 꼭 필요한 마음의 신호입니다.

기쁘고 즐거울 때는 웃음을 짓거나 "정말 신난다!"라고 말하며 기분을 나눌 수 있어요. 반대로 ( ⓐ ) 때는 억지로 참지 말고 엉엉 울어도 괜찮아요. 눈물은 마음을 씻어 내는 데 도움을 주거든요.

가장 중요한 것은 자신의 감정을 솔직하게 말하는 것입니다. 친구가 나를 때려서 화가 났다면, 같이 때리는 대신 "네가 때려서 아프고 화가 나. 하지 마."라고 말해야 합니다.
}
\begin{kfQBlock}{01}{객관식}
\kfQStem{다음 중 기분이 좋고 즐거운 감정을 나타내는 말은 무엇인가요?}
\begin{kfChoices}
\kfChoice 슬픔
\kfChoice 기쁨
\kfChoice 두려움
\kfChoice 화
\end{kfChoices}
\end{kfQBlock}

\begin{kfQBlock}{02}{객관식}
\kfQStem{다음 설명에 알맞은 감정의 이름은 무엇인가요? "갑작스러운 일을 당해서 가슴이 쿵 하고 뛰는 느낌이에요."}
\begin{kfChoices}
\kfChoice 심심함
\kfChoice 부끄러움
\kfChoice 놀람
\kfChoice 편안함
\end{kfChoices}
\end{kfQBlock}

\begin{kfQBlock}{03}{객관식}
\kfQStem{화가 났을 때의 마음을 빗대어 표현한 말은 무엇인가요?}
\begin{kfChoices}
\kfChoice 해
\kfChoice 비
\kfChoice 천둥
\kfChoice 밤
\end{kfChoices}
\end{kfQBlock}

\begin{kfQBlock}{04}{객관식}
\kfQStem{다음 중 이 동시에서 다루지 \underline{않은} 감정은 무엇인가요?}
\begin{kfChoices}
\kfChoice 기쁨
\kfChoice 슬픔
\kfChoice 부러움
\kfChoice 무서움
\end{kfChoices}
\end{kfQBlock}

\begin{kfQBlock}{05}{객관식}
\kfQStem{빈칸 ⓐ에 들어갈 알맞은 감정은 무엇인가요?}
\begin{kfChoices}
\kfChoice 기쁠
\kfChoice 슬플
\kfChoice 즐거울
\kfChoice 신날
\end{kfChoices}
\end{kfQBlock}

\begin{kfQBlock}{06}{객관식}
\kfQStem{친구가 나를 때려서 화가 났을 때 가장 올바른 행동은 무엇인가요?}
\begin{kfChoices}
\kfChoice 아무 말도 하지 않는다.
\kfChoice 친구에게 나쁜 말을 한다.
\kfChoice 솔직하게 말로 표현한다.
\kfChoice 똑같이 때려준다.
\end{kfChoices}
\end{kfQBlock}

\begin{kfQBlock}{07}{객관식}
\kfQStem{사람의 감정에 대한 설명으로 맞는 것은 무엇인가요?}
\begin{kfChoices}
\kfChoice 기쁨만 소중하고 슬픔은 쓸모없다.
\kfChoice 감정은 날씨처럼 매일 변할 수 있다.
\kfChoice 화가 날 때는 무조건 참아야 한다.
\kfChoice 어른들은 두려움을 느끼지 않는다.
\end{kfChoices}
\end{kfQBlock}

\begin{kfQBlock}{08}{서술형}
\kfQStem{이 동시에서 '무서울 때' 마음에 찾아오는 것은 무엇인지 찾아 쓰세요.}
\kfAnswerNote{답안}{4}
\end{kfQBlock}

\begin{kfQBlock}{09}{서술형}
\kfQStem{내 마음을 잘 표현하는 것이 중요한 이유는 무엇인가요? 글의 내용을 바탕으로 한 문장으로 쓰세요.}
\kfAnswerNote{답안}{4}
\end{kfQBlock}

\begin{kfQBlock}{10}{서술형}
\kfQStem{다음 문장에서 '누가'에 해당하는 말을 찾아 쓰세요.}
\begin{kfEvidence}"아기가 방긋방긋 웃어요."\end{kfEvidence}
\kfAnswerNote{답안}{4}
\end{kfQBlock}




\clearpage

\kfChapterCover{3}{소쉬르1 (챕터3)}{Chapter 3}{이번 챕터의 학습 내용을 시작합니다.}
\begin{kfChapterAreaList}
\kfChapterAreaItem{어휘}{kfAreaVocab}{핵심 어휘 익히기}
\kfChapterAreaItem{개념}{kfAreaConcept}{배경지식 넓히기: 하루를 나누는 기준}
\kfChapterAreaItem{문학}{kfAreaLiterature}{「"지수야, 일어날 시간이야!"」}
\kfChapterAreaItem{비문학}{kfAreaNonlit}{시각을 읽는 법}
\kfChapterAreaItem{실력확인}{kfAreaWeekly}{}
\end{kfChapterAreaList}

\kfStartVocab{핵심 어휘 익히기}


\begin{kfVocab}
\kfVocabItem{1}{날마다 규칙적으로 하는 일이나 하루 동안의 계획}
\kfVocabItem{2}{해가 떠서 낮 12시(정오)가 되기 전까지의 시간}
\kfVocabItem{3}{낮 12시, 해가 가장 높이 떠 있는 시각}
\kfVocabItem{4}{낮 12시(정오)부터 밤 12시(자정)까지의 시간}
\kfVocabItem{5}{정해진 기준에 따라 차례차례 이어지는 줄}
\kfVocabItem{6}{시간의 흐름 속에 있는 어느 한 시점 (예: 3시 15분)}
\end{kfVocab}


\kfStartConcept{배경지식 넓히기: 하루를 나누는 기준}


\kfPassageNoteNT{%
우리는 하루를 크게 오전과 오후로 나누어 부릅니다. 이 둘을 나누는 기준은 바로 낮 12시인 '정오'입니다.

해가 뜨고 정오가 되기 전까지를 '오전'이라고 해요. 오전에는 보통 잠에서 깨어 학교에 가거나 공부를 하며 하루를 시작하지요. 정오가 지나고 해가 지기 시작하여 밤 12시가 될 때까지는 '오후'라고 합니다. 오후에는 학교 수업을 마치고 집으로 돌아와 가족들과 시간을 보냅니다.

이렇게 시간의 순서에 따라 규칙적인 일과를 보내면 보람찬 하루를 완성할 수 있습니다.
\par\medskip\begin{itemize}[leftmargin=18pt, itemsep=2pt, topsep=2pt, label=\textbullet]
\item 오전 9시 — 학교 가는 시간
\item 정오(12시) — 점심 먹는 시간
\item 오후 3시 — 학교 마치는 시간
\end{itemize}
}

\kfPassageNoteNT{%
시간의 흐름을 나타내는 말 시간은 눈에 보이지 않지만 멈추지 않고 계속 흘러가요. 이런 시간의 순서를 나타내는 여러 가지 말이 있답니다. * 하루의 순서: 아침(해가 뜰 때)  점심(해가 높이 떴을 때)  저녁(해가 질 때)  밤(깜깜할 때) * 날짜의 순서: 어제(지나간 날)  오늘(지금의 날)  내일(다가올 날) * 계절의 순서: 봄(싹이 나요)  여름(초록잎이 무성해요)  가을(낙엽이 져요)  겨울(눈이 내려요) 이런 말들을 알면 어떤 일이 언제 일어났는지 쉽게 알 수 있어요.
\par\medskip\begin{itemize}[leftmargin=18pt, itemsep=2pt, topsep=2pt, label=\textbullet]
\item 하루: 아침 → 점심 → 저녁 → 밤
\item 날짜: 어제 → 오늘 → 내일
\item 계절: 봄 → 여름 → 가을 → 겨울
\end{itemize}
}

\kfActivityBg \par\kfMaybeNeedspace{24\baselineskip}\noindent{\kfTipMark\,\,}{\small\color{kfInk} 다음 단어와 알맞은 설명을 선으로 연결해 보세요.}\par\nopagebreak[4]\smallskip\nopagebreak[4]
\begin{kfConnect}
\kfPair{일과}{낮 12시(정오)부터 밤까지의 시간이에요}
\kfPair{오전}{하루 동안 해야 할 규칙적인 일이에요}
\kfPair{정오}{낮 12시가 되기 전까지의 시간이에요}
\kfPair{오후}{시간의 어느 한 시점을 말해요}
\kfPair{시각}{낮 12시를 가리키는 말이에요}
\end{kfConnect}

\kfActivityBg \par\kfMaybeNeedspace{24\baselineskip}\noindent{\kfTipMark\,\,}{\small\color{kfInk} 다음 문장이 옳으면 O, 옳지 않으면 X를 표시하세요.}\par\nopagebreak[4]\smallskip\nopagebreak[4]
\begin{kfOXList}
\kfOXItem 하루를 오전과 오후로 나누는 기준은 '정오'입니다. \kfFillBlank[12mm]{} \kfOX
\kfOXItem 오후는 해가 뜨기 시작해서 낮 12시가 될 때까지를 말합니다. \kfFillBlank[12mm]{} \kfOX
\kfOXItem '일과'는 날마다 규칙적으로 하는 하루 동안의 일을 뜻합니다. \kfFillBlank[12mm]{} \kfOX
\end{kfOXList}

\kfActivityBg \par\kfMaybeNeedspace{24\baselineskip}\noindent{\kfTipMark\,\,}{\small\color{kfInk} 다음 단어들을 시간의 흐름 순서대로 나열하여 번호를 쓰세요.}\par\nopagebreak[4]\smallskip\nopagebreak[4]
\begin{kfBlankList}
\kfBlankItem \kfFillBlank[20mm]{} 정오 \kfFillBlank[20mm]{} 오전 \kfFillBlank[20mm]{} 오후
\kfBlankItem \kfFillBlank[20mm]{} 내일 \kfFillBlank[20mm]{} 어제 \kfFillBlank[20mm]{} 오늘
\kfBlankItem \kfFillBlank[20mm]{} 점심 \kfFillBlank[20mm]{} 아침 \kfFillBlank[20mm]{} 밤 \kfFillBlank[20mm]{} 저녁
\kfBlankItem \kfFillBlank[20mm]{} 여름 \kfFillBlank[20mm]{} 봄 \kfFillBlank[20mm]{} 겨울 \kfFillBlank[20mm]{} 가을
\end{kfBlankList}


\kfStartLit{「"지수야, 일어날 시간이야!"」}


\kfPassageNote{}{%
"지수야, 일어날 시간이야!"

오전 8시, 지수는 알람 소리에 눈을 떴습니다. 오늘은 학교에서 현장 학습을 가는 날이라 평소보다 일찍 일과를 시작했습니다. 지수는 세수를 하고 든든하게 아침밥을 먹은 뒤 집을 나섰습니다.

정오가 가까워지자 친구들과 함께 박물관 앞 잔디밭에 모였습니다. "지금 시각은 12시 정각입니다. 점심 식사를 하세요!" 선생님의 말씀에 따라 지수는 친구들과 도시락을 먹으며 즐거운 시간을 보냈습니다.

현장 학습을 마치고 돌아온 오후 4시, 지수는 태권도 학원에 갔습니다. 힘차게 발차기 연습을 하고 나니 땀이 흘렀지만 기분은 상쾌했습니다. 집에 돌아와 씻고 나니 어느덧 해가 지고 어두워졌습니다. 지수는 일기장에 오늘 있었던 일을 순서대로 적으며 하루를 마무리했습니다.
}


\kfActivityBg \par\kfMaybeNeedspace{18\baselineskip}\noindent{\kfTipMark\,\,}{\small\color{kfInk} 주어진 안내에 맞춰 자유롭게 글로 표현해 보세요.}\par\nopagebreak[4]\smallskip\nopagebreak[4]
\kfWritingPrompt{내일 꼭 지키고 싶은 일을 '오전'과 '오후'로 나누어 계획해 보세요.}
\kfWriting{답안 작성}{8}

\kfSecHeader{kfAreaLiterature}{문제}{}

\begin{kfQBlock}{01}{객관식}
\kfQStem{지수가 오전 8시에 일어나서 가장 먼저 한 일은 무엇인가요?}
\begin{kfChoices}
\kfChoice 도시락 먹기
\kfChoice 태권도 연습하기
\kfChoice 세수 하기
\kfChoice 일기 쓰기
\end{kfChoices}
\end{kfQBlock}
\begin{kfQBlock}{02}{객관식}
\kfQStem{지수와 친구들이 점심 도시락을 먹은 때는 언제인가요?}
\begin{kfChoices}
\kfChoice 오전 9시
\kfChoice 정오
\kfChoice 오후 4시
\kfChoice 밤 10시
\end{kfChoices}
\end{kfQBlock}
\begin{kfQBlock}{03}{객관식}
\kfQStem{다음 중 지수의 오후 일과에 해당하는 것은 무엇인가요?}
\begin{kfChoices}
\kfChoice 늦잠 자기
\kfChoice 태권도 학원 가기
\kfChoice 아침밥 먹기
\kfChoice 박물관 도착하기
\end{kfChoices}
\end{kfQBlock}
\begin{kfQBlock}{04}{서술형}
\kfQStem{여러분의 하루 일과 중 가장 중요하게 생각하는 일은 무엇인가요?}
\kfAnswerNote{답안}{4}
\end{kfQBlock}
\par\noindent\textit{\small * 나는 \kfFillBlank[12mm]{}을/를 가장 중요하게 생각해요. * 왜냐하면 \kfFillBlank[12mm]{}}


\kfStartNonlit{시각을 읽는 법}


\kfPassageNote{시각을 읽는 법}{%
우리는 정확한 시각을 알기 위해 시계를 봅니다. 시계에는 길이가 다른 두 개의 바늘이 부지런히 움직이고 있습니다. 이 바늘들은 각각 '시'와 '분'을 가리키며 우리가 규칙적인 일과를 보낼 수 있도록 도와줍니다.

짧고 굵은 바늘은 '시'를 나타내는 시침입니다. 시침이 숫자 9를 가리키면 9시, 10을 가리키면 10시라고 읽습니다. 시침은 아주 느리게 움직여서, 오전에 한 바퀴, 오후에 한 바퀴, 하루에 총 두 바퀴를 돕니다.

길고 얇은 바늘은 '분'을 나타내는 분침입니다. 분침이 숫자 12를 가리키면 '0분'이 되어 정각을 의미합니다. 예를 들어, 시침이 3에 있고 분침이 12에 있으면 '3시 정각'이 됩니다. 시침과 분침이 모두 12를 가리키는 정오가 되면 점심을 먹고, 밤 12시가 되면 하루가 끝납니다.
}


\kfActivityBg \par\needspace{15\baselineskip}
\kfSecHeader{kfAreaNonlit}{활동}{문장 독해}
\par\kfMaybeNeedspace{32\baselineskip}\noindent{\kfTipMark\,\,}{\small\color{kfInk} 각 문장을 읽고 물음에 답하세요.}\par\nopagebreak[4]\smallskip\nopagebreak[4]
\begin{kfSentenceUnit}{1}{나는 오전에 책을}
\kfSentQ{1.}{(1) 이 문장에서 '무엇을'에 해당하는 것은 무엇인가요?}
\begin{kfChoices}
\kfChoice 나는
\kfChoice 오전에
\kfChoice 책을
\end{kfChoices}
\kfSentQ{1.}{(2) 이 문장에서 '어찌하다'에 해당하는 것은 무엇인가요?}
\begin{kfChoices}
\kfChoice 나는
\kfChoice 오전에
\kfChoice 읽었어요
\end{kfChoices}
\end{kfSentenceUnit}

\begin{kfSentenceUnit}{2}{우리는 오후에 그림을}
\kfSentQ{2.}{(1) 이 문장에서 '무엇을'에 해당하는 것은 무엇인가요?}
\begin{kfChoices}
\kfChoice 우리는
\kfChoice 오후에
\kfChoice 그림을
\end{kfChoices}
\kfSentQ{2.}{(2) 이 문장에서 '어찌하다'에 해당하는 것은 무엇인가요?}
\begin{kfChoices}
\kfChoice 우리는
\kfChoice 오후에
\kfChoice 그렸어요
\end{kfChoices}
\end{kfSentenceUnit}

\kfActivityBg \par\kfMaybeNeedspace{24\baselineskip}\noindent{\kfTipMark\,\,}{\small\color{kfInk} 빈칸에 알맞은 말을 채워 보세요.}\par\nopagebreak[4]\smallskip\nopagebreak[4]
\begin{kfBlankList}
\kfBlankItem 빈칸에 알맞은 말을 <보기>에서 찾아 써 보세요.

* 시계의 짧은 바늘은 \kfFillBlank[42mm]{}을/를 알려 줍니다. * 시계의 긴 바늘은 \kfFillBlank[42mm]{}을/를 알려 줍니다. * 시침과 분침이 모두 12를 가리키는 낮 시간을 \kfFillBlank[42mm]{}라고 합니다.
\begin{kfEvidence}시각, 시, 분, 정오, 오전, 오후\end{kfEvidence}
\end{kfBlankList}

\kfActivityBg \par\kfMaybeNeedspace{18\baselineskip}\noindent{\kfTipMark\,\,}{\small\color{kfInk} 주어진 안내에 맞춰 자유롭게 글로 표현해 보세요.}\par\nopagebreak[4]\smallskip\nopagebreak[4]
\kfWritingPrompt{하루 중에 내가 가장 좋아하는 시간은 언제인가요? 무엇을 할 때인지 써 보세요.}
\kfWriting{답안 작성}{8}

\kfSecHeader{kfAreaNonlit}{문제}{}

\begin{kfQBlock}{01}{객관식}
\kfQStem{시계의 짧은 바늘(시침)에 대한 설명으로 옳은 것은 무엇인가요?}
\begin{kfChoices}
\kfChoice '분'을 알려 준다.
\kfChoice 하루에 한 바퀴만 돈다.
\kfChoice 아주 빨리 움직인다.
\kfChoice '시'를 알려 준다.
\end{kfChoices}
\end{kfQBlock}
\begin{kfQBlock}{02}{객관식}
\kfQStem{'정오'는 시계 바늘이 어떻게 되어 있을 때인가요?}
\begin{kfChoices}
\kfChoice 짧은 바늘은 6, 긴 바늘은 12
\kfChoice 짧은 바늘은 12, 긴 바늘은 6
\kfChoice 짧은 바늘과 긴 바늘 모두 12
\kfChoice 짧은 바늘과 긴 바늘 모두 3
\end{kfChoices}
\end{kfQBlock}
\begin{kfQBlock}{03}{객관식}
\kfQStem{긴 바늘(분침)이 숫자 12를 가리킬 때 분은 어떻게 읽나요?}
\begin{kfChoices}
\kfChoice 12분
\kfChoice 30분
\kfChoice 60분
\kfChoice 0분 (정각)
\end{kfChoices}
\end{kfQBlock}


\kfStartWeekly{실력확인 영역 학습}


\noindent\textit{\small 제한시간: 20분 / 총 10문항}\par\medskip
\par\noindent\textbf{[3\textasciitilde{}5] 다음 글을 읽고 물음에 답하시오.}\par\medskip
\kfPassageNote{지문 3}{%
지수의 알찬 일과

"지수야, 일어날 시간이야!"

( ㉠ ) 8시, 지수는 알람 소리에 눈을 떴습니다. 오늘은 현장 학습을 가는 날이라 일찍 하루를 시작했습니다.

정오가 가까워지자 친구들과 함께 박물관 앞 잔디밭에 모였습니다. "지금 시각은 12시 정각입니다. 식사를 하세요!" 선생님의 말씀에 따라 지수는 친구들과 도시락을 먹었습니다.

현장 학습을 마치고 돌아온 오후 4시, 지수는 태권도 학원에 갔습니다. 집에 돌아와 씻고 나니 어느덧 해가 지고 어두워졌습니다. 지수는 일기장에 오늘 있었던 일을 순서대로 적으며 하루를 마무리했습니다.
}
\par\noindent\textbf{[6\textasciitilde{}8] 다음 글을 읽고 물음에 답하시오.}\par\medskip
\kfPassageNote{지문 6}{%
시각을 읽는 법

우리는 정확한 시각을 알기 위해 시계를 봅니다. 시계에는 길이가 다른 두 개의 바늘이 있습니다.

짧고 굵은 바늘은 '시'를 나타내는 시침입니다. 시침이 숫자 9를 가리키면 9시라고 읽습니다. 시침은 아주 느리게 움직여서 하루에 두 바퀴를 돕니다.

길고 얇은 바늘은 '분'을 나타내는 분침입니다. 분침이 숫자 12를 가리키면 '0분'이 되어 정각을 의미합니다. 예를 들어, 짧은 바늘이 3에 있고 긴 바늘이 12에 있으면 ( ⓐ ) 정각이 됩니다.
}
\begin{kfQBlock}{01}{객관식}
\kfQStem{해가 떠서 낮 12시가 되기 전까지의 시간을 무엇이라고 하나요?}
\begin{kfChoices}
\kfChoice 오전
\kfChoice 정오
\kfChoice 오후
\kfChoice 자정
\end{kfChoices}
\end{kfQBlock}

\begin{kfQBlock}{02}{객관식}
\kfQStem{날마다 규칙적으로 하는 하루의 일을 뜻하는 단어는 무엇인가요?}
\begin{kfChoices}
\kfChoice 시각
\kfChoice 순서
\kfChoice 일과
\kfChoice 연습
\end{kfChoices}
\end{kfQBlock}

\begin{kfQBlock}{03}{객관식}
\kfQStem{이 글에서 지수가 '정오'에 한 일은 무엇인가요?}
\begin{kfChoices}
\kfChoice 세수 하기
\kfChoice 도시락 먹기
\kfChoice 태권도 하기
\kfChoice 일기 쓰기
\end{kfChoices}
\end{kfQBlock}

\begin{kfQBlock}{04}{객관식}
\kfQStem{지수가 하루를 마무리하며 한 행동은 무엇인가요?}
\begin{kfChoices}
\kfChoice 알람 끄기
\kfChoice 박물관 견학
\kfChoice 도시락 먹기
\kfChoice 일기장에 일 적기
\end{kfChoices}
\end{kfQBlock}

\begin{kfQBlock}{05}{객관식}
\kfQStem{시계의 긴 바늘(분침)이 12를 가리킬 때의 설명으로 맞는 것은?}
\begin{kfChoices}
\kfChoice 12분이라고 읽는다.
\kfChoice 30분을 나타낸다.
\kfChoice 0분인 정각을 의미한다.
\kfChoice 시침과 똑같은 역할을 한다.
\end{kfChoices}
\end{kfQBlock}

\begin{kfQBlock}{06}{객관식}
\kfQStem{빈칸 ⓐ에 들어갈 알맞은 시각은 무엇인가요?}
\begin{kfChoices}
\kfChoice 3시
\kfChoice 12시
\kfChoice 3시 12분
\kfChoice 12시 3분
\end{kfChoices}
\end{kfQBlock}

\begin{kfQBlock}{07}{객관식}
\kfQStem{다음 중 시간의 흐름 순서가 올바른 것은 무엇인가요?}
\begin{kfChoices}
\kfChoice 오후 → 정오 → 오전
\kfChoice 정오 → 오전 → 오후
\kfChoice 오전 → 정오 → 오후
\kfChoice 오후 → 오전 → 정오
\end{kfChoices}
\end{kfQBlock}

\begin{kfQBlock}{08}{서술형}
\kfQStem{빈칸 ㉠에 들어갈 알맞은 말을 '오전' 또는 '오후' 중에서 골라 쓰세요.}
\kfAnswerNote{답안}{4}
\end{kfQBlock}

\begin{kfQBlock}{09}{서술형}
\kfQStem{시계의 짧은 바늘(시침)이 하는 역할은 무엇인가요? 글의 내용을 바탕으로 쓰세요.}
\kfAnswerNote{답안}{4}
\end{kfQBlock}

\begin{kfQBlock}{10}{서술형}
\kfQStem{다음 문장에서 '무엇을'에 해당하는 말을 찾아 쓰세요.}
\begin{kfEvidence}"지수는 태권도 학원에서 발차기를 연습했어요."\end{kfEvidence}
\kfAnswerNote{답안}{4}
\end{kfQBlock}




\clearpage

\kfChapterCover{4}{소쉬르1 (챕터4)}{Chapter 4}{이번 챕터의 학습 내용을 시작합니다.}
\begin{kfChapterAreaList}
\kfChapterAreaItem{어휘}{kfAreaVocab}{핵심 어휘 익히기}
\kfChapterAreaItem{개념}{kfAreaConcept}{배경지식 넓히기: 마음을 전하는 '편지'}
\kfChapterAreaItem{문학}{kfAreaLiterature}{「할머니, 안녕하세요? 저 민수예요.」}
\kfChapterAreaItem{비문학}{kfAreaNonlit}{우리 가족을 소개합니다}
\kfChapterAreaItem{실력확인}{kfAreaWeekly}{}
\end{kfChapterAreaList}

\kfStartVocab{핵심 어휘 익히기}


\begin{kfVocab}
\kfVocabItem{1}{서로 뜻이 잘 맞고 정다워서 분위기가 좋음}
\kfVocabItem{2}{정성을 다해 보호하고 돕는 일}
\kfVocabItem{3}{편안하게 잘 지내는지에 대한 소식}
\kfVocabItem{4}{용기나 의욕이 솟아나도록 북돋아 줌}
\kfVocabItem{5}{윗어른을 공손히 받들어 모심}
\kfVocabItem{6}{형제나 자매 사이의 깊은 정과 사랑}
\end{kfVocab}


\kfStartConcept{배경지식 넓히기: 마음을 전하는 '편지'}


\kfPassageNoteNT{%
가족이나 친구에게 말로 하기 부끄러운 마음을 전하고 싶을 때 우리는 '편지'를 씁니다. 편지는 정해진 형식이 있어요.

가장 먼저 '받는 사람'을 쓰고, 날씨나 건강을 묻는 '첫인사'를 합니다. 그다음에는 하고 싶은 말을 자세히 쓰고, 마지막으로 행복을 비는 '끝인사'를 남깁니다. 그리고 편지를 쓴 날짜와 '보내는 사람'의 이름을 적습니다.

멀리 떨어져 있어도 안부를 묻고 서로를 격려하는 편지는 가족의 사랑을 더욱 깊어지게 만들어 줍니다.
\par\medskip\begin{itemize}[leftmargin=18pt, itemsep=2pt, topsep=2pt, label=\textbullet]
\item 받는 사람: '사랑하는 할머니께'
\item 첫인사: '잘 지내고 계신가요?'
\item 끝인사: '항상 건강하세요.'
\item 보내는 사람: '○○○ 올림'
\end{itemize}
}

\kfActivityBg \par\kfMaybeNeedspace{24\baselineskip}\noindent{\kfTipMark\,\,}{\small\color{kfInk} 다음 단어와 알맞은 설명을 선으로 연결해 보세요.}\par\nopagebreak[4]\smallskip\nopagebreak[4]
\begin{kfConnect}
\kfPair{화목하다}{윗어른을 공손히 모시는 마음이에요}
\kfPair{보살핌}{용기를 낼 수 있도록 응원해 줘요}
\kfPair{안부}{사이가 좋고 정다운 모습이에요}
\kfPair{격려}{잘 지내는지 묻는 소식이에요}
\kfPair{공경}{정성을 다해 보호하고 돌봐 줘요}
\end{kfConnect}

\kfActivityBg \par\kfMaybeNeedspace{24\baselineskip}\noindent{\kfTipMark\,\,}{\small\color{kfInk} 다음 문장이 옳으면 O, 옳지 않으면 X를 표시하세요.}\par\nopagebreak[4]\smallskip\nopagebreak[4]
\begin{kfOXList}
\kfOXItem 편지는 받는 사람, 첫인사, 내용, 끝인사, 보낸 사람의 순서로 씁니다. \kfFillBlank[12mm]{} \kfOX
\kfOXItem 편지는 말로 하기 힘든 마음을 전할 때 좋은 방법입니다. \kfFillBlank[12mm]{} \kfOX
\kfOXItem '안부'는 윗어른을 공손하게 모신다는 뜻입니다. \kfFillBlank[12mm]{} \kfOX
\end{kfOXList}

\kfActivityBg \par\kfMaybeNeedspace{24\baselineskip}\noindent{\kfTipMark\,\,}{\small\color{kfInk} 편지를 쓸 때 들어갈 내용으로 알맞은 곳에 번호를 쓰세요.}\par\nopagebreak[4]\smallskip\nopagebreak[4]
\begin{kfBlankList}
\kfBlankItem ( 1 ) 받는 사람 ( 2 ) 첫인사 ( 3 ) 쓴 날짜와 보낸 사람

* \kfFillBlank[20mm]{} 사랑하는 할머니께 * \kfFillBlank[20mm]{} 2024년 5월 8일 손자 민수 올림 * \kfFillBlank[20mm]{} 할머니, 그동안 건강하셨나요?
\end{kfBlankList}


\kfStartLit{「할머니, 안녕하세요? 저 민수예요.」}


\kfPassageNote{}{%
할머니, 안녕하세요? 저 민수예요.

시골에는 하얀 눈이 많이 내렸나요? 이곳은 날씨가 많이 추워졌어요. 감기에 걸리시지는 않았는지 걱정이 돼요.

저는 요즘 태권도를 열심히 배우고 있어요. 처음에는 발차기가 어려웠는데, 관장님의 격려 덕분에 지금은 아주 잘해요. 이번 설날에 할머니 댁에 가면 멋진 발차기 시범을 보여 드릴게요.

그리고 제가 좋아하는 붕어빵도 사서 갈게요. 할머니께서 끓여 주시는 따뜻한 된장찌개가 정말 먹고 싶어요. 할머니, 우리 가족 모두 할머니를 많이 보고 싶어 해요.

설날까지 건강하게 지내세요. 제가 가서 어깨 안마를 시원하게 해 드릴게요.

안녕히 계세요.

2024년 12월 20일

할머니를 사랑하는 손자 민수 올림
}


\kfActivityBg \par\kfMaybeNeedspace{18\baselineskip}\noindent{\kfTipMark\,\,}{\small\color{kfInk} 주어진 안내에 맞춰 자유롭게 글로 표현해 보세요.}\par\nopagebreak[4]\smallskip\nopagebreak[4]
\kfWritingPrompt{여러분이 사랑하는 가족 중 한 명에게 하고 싶은 말을 짧게 써 보세요.}
\kfWriting{답안 작성}{8}

\kfSecHeader{kfAreaLiterature}{문제}{}

\begin{kfQBlock}{01}{객관식}
\kfQStem{이 편지를 받은 사람은 누구인가요?}
\begin{kfChoices}
\kfChoice 민수
\kfChoice 관장님
\kfChoice 할머니
\kfChoice 어머니
\end{kfChoices}
\end{kfQBlock}
\begin{kfQBlock}{02}{객관식}
\kfQStem{민수가 요즘 열심히 배우고 있는 것은 무엇인가요?}
\begin{kfChoices}
\kfChoice 피아노
\kfChoice 태권도
\kfChoice 수영
\kfChoice 미술
\end{kfChoices}
\end{kfQBlock}
\begin{kfQBlock}{03}{객관식}
\kfQStem{민수가 설날에 할머니 댁에 가서 하겠다고 약속한 일이 \underline{아닌} 것은 무엇인가요?}
\begin{kfChoices}
\kfChoice 발차기 시범 보이기
\kfChoice 붕어빵 사 가기
\kfChoice 어깨 안마 해 드리기
\kfChoice 된장찌개 끓여 드리기
\end{kfChoices}
\end{kfQBlock}
\begin{kfQBlock}{04}{서술형}
\kfQStem{민수는 할머니를 생각하며 어떤 마음으로 편지를 썼을까요?}
\kfAnswerNote{답안}{4}
\end{kfQBlock}
\par\noindent\textit{\small * 민수는 할머니가 \kfFillBlank[12mm]{} 마음으로 편지를 썼을 거예요.}


\kfStartNonlit{우리 가족을 소개합니다}


\kfPassageNote{우리 가족을 소개합니다}{%
우리 가족은 모두 네 명입니다. 아빠, 엄마, 형, 그리고 저입니다. 우리 가족은 항상 서로를 아끼고 보살핌을 주고받는 화목한 가족입니다.

아빠는 주말마다 우리와 함께 놀아 주는 것을 좋아하십니다. 특히 등산을 좋아하셔서 우리 가족은 자주 산에 갑니다. 아빠는 힘들어하는 저를 엎어 주기도 하시고 "조금만 더 힘내자!"라고 격려해 주십니다.

엄마는 요리 솜씨가 아주 좋으십니다. 우리 가족이 좋아하는 음식을 뚝딱 만들어 주십니다. 엄마는 항상 웃는 얼굴로 우리 이야기를 잘 들어 주셔서, 저는 학교에서 있었던 일을 엄마에게 제일 먼저 이야기합니다.

형은 장난꾸러기지만 저를 잘 챙겨 줍니다. 제가 숙제를 어려워하면 친절하게 가르쳐 주고, 심심할 때는 재미있는 놀이를 함께 합니다. 우리는 우애가 깊은 형제입니다. 저는 우리 가족을 정말 사랑합니다.
}


\kfActivityBg \par\needspace{15\baselineskip}
\kfSecHeader{kfAreaNonlit}{활동}{문장 독해}
\par\kfMaybeNeedspace{32\baselineskip}\noindent{\kfTipMark\,\,}{\small\color{kfInk} 각 문장을 읽고 물음에 답하세요.}\par\nopagebreak[4]\smallskip\nopagebreak[4]
\begin{kfSentenceUnit}{1}{민수는 할머니를 사랑해요}
\kfSentQ{1.}{(1) 이 문장에서 '누가'에 해당하는 것은 무엇인가요?}
\begin{kfChoices}
\kfChoice 민수는
\kfChoice 할머니를
\kfChoice 사랑해요
\end{kfChoices}
\kfSentQ{1.}{(2) 이 문장에서 '어찌하다'에 해당하는 것은 무엇인가요?}
\begin{kfChoices}
\kfChoice 민수는
\kfChoice 할머니를
\kfChoice 사랑해요
\end{kfChoices}
\end{kfSentenceUnit}

\begin{kfSentenceUnit}{2}{아빠는 저를 업어}
\kfSentQ{2.}{(1) 이 문장에서 '누가'에 해당하는 것은 무엇인가요?}
\begin{kfChoices}
\kfChoice 아빠는
\kfChoice 저를
\kfChoice 업어
\end{kfChoices}
\kfSentQ{2.}{(2) 이 문장에서 '누구를'에 해당하는 것은 무엇인가요?}
\begin{kfChoices}
\kfChoice 아빠는
\kfChoice 저를
\kfChoice 업어
\end{kfChoices}
\end{kfSentenceUnit}

\kfActivityBg \par\kfMaybeNeedspace{24\baselineskip}\noindent{\kfTipMark\,\,}{\small\color{kfInk} 빈칸에 알맞은 말을 채워 보세요.}\par\nopagebreak[4]\smallskip\nopagebreak[4]
\begin{kfBlankList}
\kfBlankItem 빈칸에 알맞은 말을 <보기>에서 찾아 써 보세요.

* 아빠는 등산을 할 때 힘들어하는 저를 \kfFillBlank[42mm]{}해 주십니다. * 엄마는 요리 솜씨가 좋고, 제 이야기를 잘 \kfFillBlank[42mm]{} 주십니다. * 형은 장난꾸러기지만 저를 잘 챙겨 주는 \kfFillBlank[42mm]{}가 깊은 형입니다.
\begin{kfEvidence}우애, 들어, 격려, 화목, 친절\end{kfEvidence}
\end{kfBlankList}

\kfActivityBg \par\kfMaybeNeedspace{18\baselineskip}\noindent{\kfTipMark\,\,}{\small\color{kfInk} 주어진 안내에 맞춰 자유롭게 글로 표현해 보세요.}\par\nopagebreak[4]\smallskip\nopagebreak[4]
\kfWritingPrompt{우리 가족 중 한 사람을 골라 자랑하고 싶은 점을 써 보세요.}
\kfWriting{답안 작성}{8}

\kfSecHeader{kfAreaNonlit}{문제}{}

\begin{kfQBlock}{01}{객관식}
\kfQStem{아빠가 주말에 가족들과 함께 즐겨 하시는 일은 무엇인가요?}
\begin{kfChoices}
\kfChoice 요리하기
\kfChoice 등산하기
\kfChoice 청소하기
\kfChoice 낮잠 자기
\end{kfChoices}
\end{kfQBlock}
\begin{kfQBlock}{02}{객관식}
\kfQStem{글쓴이가 형을 소개한 내용으로 알맞은 것은 무엇인가요?}
\begin{kfChoices}
\kfChoice 한 번도 장난을 친 적이 없는 점잖은 형이다.
\kfChoice 글쓴이가 심심하면 학원에 데려다주신다.
\kfChoice 글쓴이의 숙제를 대신 풀어서 끝내 준다.
\kfChoice 숙제가 어려울 때 친절하게 가르쳐 준다.
\end{kfChoices}
\end{kfQBlock}
\begin{kfQBlock}{03}{객관식}
\kfQStem{글쓴이가 학교에서 있었던 일을 엄마에게 제일 먼저 이야기하는 이유는 무엇인가요?}
\begin{kfChoices}
\kfChoice 가족 중 요리 솜씨가 가장 좋아서
\kfChoice 웃는 얼굴로 이야기를 잘 들어 주셔서
\kfChoice 주말마다 함께 등산을 가 주셔서
\kfChoice 어려운 숙제를 친절히 가르쳐 주셔서
\end{kfChoices}
\end{kfQBlock}


\kfStartWeekly{실력확인 영역 학습}


\noindent\textit{\small 제한시간: 20분 / 총 10문항}\par\medskip
\par\noindent\textbf{[3\textasciitilde{}5] 다음 글을 읽고 물음에 답하시오.}\par\medskip
\kfPassageNote{지문 3}{%
사랑하는 할머니께 \\[14pt]
( ㉠ ), 안녕하세요? 저 민수예요. \\[14pt]
이곳은 날씨가 많이 추워졌어요. 감기에 걸리시지는 않았는지 걱정이 돼요. \\[14pt]
저는 요즘 태권도를 열심히 배우고 있어요. 처음에는 어려웠지만 관장님의 격려 덕분에 지금은 아주 잘해요. 이번 설날에 할머니 댁에 가면 멋진 발차기 시범을 보여 드릴게요. \\[14pt]
설날까지 건강하게 지내세요. 제가 가서 어깨 안마를 시원하게 해 드릴게요. \\[14pt]
안녕히 계세요. \\[14pt]
2024년 12월 20일 \\[14pt]
할머니를 사랑하는 손자 민수 올림
}
\par\noindent\textbf{[6\textasciitilde{}8] 다음 글을 읽고 물음에 답하시오.}\par\medskip
\kfPassageNote{지문 6}{%
우리 가족을 소개합니다

우리 가족은 아빠, 엄마, 형, 그리고 저까지 모두 네 명입니다. 우리는 서로를 아끼는 화목한 가족입니다.

아빠는 주말마다 우리와 등산하는 것을 좋아하십니다. 아빠는 힘들어하는 저를 엎어 주기도 하시고 "힘내자!"라고 응원해 주십니다.

엄마는 ( ⓐ ) 솜씨가 아주 좋으십니다. 우리 가족이 좋아하는 음식을 뚝딱 만들어 주시고, 제 이야기를 항상 잘 들어 주십니다.

형은 장난꾸러기지만 저를 잘 챙겨 줍니다. 제가 숙제를 어려워하면 친절하게 가르쳐 줍니다.
}
\begin{kfQBlock}{01}{객관식}
\kfQStem{형제나 자매 사이의 깊은 정과 사랑을 뜻하는 말은 무엇인가요?}
\begin{kfChoices}
\kfChoice 공경
\kfChoice 우애
\kfChoice 안부
\kfChoice 격려
\end{kfChoices}
\end{kfQBlock}

\begin{kfQBlock}{02}{객관식}
\kfQStem{다음 중 '화목하다'라는 말과 가장 잘 어울리는 가족의 모습은 무엇인가요?}
\begin{kfChoices}
\kfChoice 서로 돕고 웃음꽃이 피는 가족
\kfChoice 매일 큰 소리로 다투는 가족
\kfChoice 서로에게 아무 관심이 없는 가족
\kfChoice 밥을 따로따로 먹는 가족
\end{kfChoices}
\end{kfQBlock}

\begin{kfQBlock}{03}{객관식}
\kfQStem{민수가 할머니께 편지를 쓴 이유로 가장 알맞은 것은 무엇인가요?}
\begin{kfChoices}
\kfChoice 용돈을 달라고 부탁하기 위해서
\kfChoice 친구와 싸워서 이르기 위해서
\kfChoice 안부를 묻고 소식을 전하기 위해서
\kfChoice 맛있는 음식을 보내 달라고 하기 위해서
\end{kfChoices}
\end{kfQBlock}

\begin{kfQBlock}{04}{객관식}
\kfQStem{윗글의 내용과 다른 것은 무엇인가요?}
\begin{kfChoices}
\kfChoice 민수는 요즘 태권도를 배우고 있다.
\kfChoice 민수는 설날에 할머니 댁에 갈 예정이다.
\kfChoice 민수는 할머니의 건강을 걱정하고 있다.
\kfChoice 민수는 할머니께 전화를 걸어 이야기했다.
\end{kfChoices}
\end{kfQBlock}

\begin{kfQBlock}{05}{객관식}
\kfQStem{아빠가 저에게 "힘내자!"라고 말해주신 행동을 뜻하는 단어는 무엇인가요?}
\begin{kfChoices}
\kfChoice 격려
\kfChoice 후회
\kfChoice 사과
\kfChoice 부탁
\end{kfChoices}
\end{kfQBlock}

\begin{kfQBlock}{06}{객관식}
\kfQStem{빈칸 ⓐ에 들어갈 알맞은 말은 무엇인가요?}
\begin{kfChoices}
\kfChoice 등산
\kfChoice 요리
\kfChoice 청소
\kfChoice 노래
\end{kfChoices}
\end{kfQBlock}

\begin{kfQBlock}{07}{객관식}
\kfQStem{편지를 쓸 때 가장 마지막에 들어가는 내용은 무엇인가요?}
\begin{kfChoices}
\kfChoice 받는 사람 이름
\kfChoice 날씨와 계절 인사
\kfChoice 쓴 날짜와 보낸 사람
\kfChoice 하고 싶은 말
\end{kfChoices}
\end{kfQBlock}

\begin{kfQBlock}{08}{서술형}
\kfQStem{빈칸 ㉠에 들어갈 편지를 받는 사람은 누구인지 쓰세요.}
\kfAnswerNote{답안}{4}
\end{kfQBlock}

\begin{kfQBlock}{09}{서술형}
\kfQStem{글쓴이가 우리 가족을 '화목하다'고 소개한 이유는 무엇일까요? 글의 내용을 바탕으로 쓰세요.}
\kfAnswerNote{답안}{4}
\end{kfQBlock}

\begin{kfQBlock}{10}{서술형}
\kfQStem{다음 문장에서 '누구를'에 해당하는 말을 찾아 쓰세요.}
\begin{kfEvidence}"형은 동생을 도와주었습니다."\end{kfEvidence}
\kfAnswerNote{답안}{4}
\end{kfQBlock}




\end{document}
